% (c) 2012 - 2014 Dimitrios Vrettos - d.vrettos@gmail.com
% (c) 2014 Claudio Carboncini - claudio.carboncini@gmail.com
% (c) 2014 Daniele Zambelli - daniele.zambelli@gmail.com

\chapter{Numeri razionali (e irrazionali)}

\section{Premessa storica}
\label{sec:03_origine}

Quando si deve dividere una certa grandezza o totalità in un certo numero di 
parti uguali non sempre sono sufficienti i numeri interi per rappresentare 
il risultato della divisione. Per esempio, per dividere l'unità in due parti 
uguali i numeri interi non sono sufficienti.

Gli antichi hanno affrontato questo tipo di problema utilizzando varie 
scritture per rappresentare le parti in cui dividere l'unità, ossia le 
frazioni.

I Babilonesi scrivevano frazioni aventi come denominatore una potenza di~60, 
la base della loro numerazione; tuttavia non usavano una notazione specifica 
per le frazioni ed il valore corretto andava interpretato dal contesto.

Gli Egizi facevano largo uso dei numeri frazionari che rappresentavano come 
somme di frazioni unitarie, ossia frazioni con numeratore uno. 
La frazione unitaria~$\frac{1}{n}$ (con~$n$ numero naturale diverso da zero)
veniva rappresentata in forma geroglifica ponendo il denominatore~$n$ scritto 
con la normale rappresentazione del numero~$n$ sotto ad un ovale. 
La frazione~$\frac{1}{12}$, per esempio, veniva così rappresentata:

% \begin{center}
%  % (c) 2012 Dimitrios Vrettos - d.vrettos@gmail.com
\begin{hieroglyph}{\leavevmode \loneSign{{\Hsmaller\Aca GD/66/}}}\end{hieroglyph}.
.
% \end{center}

\vspace{-2ex}
\begin{center}
\begin{inaccessibleblock}[Figura: TODO]
 \includegraphics[scale=0.28]{img/giero.png}
\end{inaccessibleblock}
\end{center}
\vspace{-2ex}
                                                                  
Nel \emph{papiro di Ahmes} (detto anche \emph{papiro di Rhind}) troviamo una 
tabella che dà la scomposizione in frazioni unitarie delle frazioni del 
tipo~$\frac{2}{n}$, con~$n$ dispari: la frazione~$\frac{2}{43}$
è rappresentata come somma di frazioni unitarie nel seguente modo:

\[\frac{2}{43}=\frac{1}{42}+\frac{1}{86}+\frac{1}{129}+\frac{1}{301}.\]
\begin{wrapfigure}{r}{0pt}
% (c) 2012 Dimitrios Vrettos - d.vrettos@gmail.com
% Occhio di Horus 
\pgfdeclareimage[interpolate=true]{eye}{\folder img/eye.png}%{\folder img/eye.pdf}
\begin{tikzpicture}[scale=.7]
    \pgftext[at=\pgfpoint {0mm}{0mm},left,base]{\pgfuseimage{eye}}
    \node (a) at (10mm,65mm) {$\frac{1}{2}$};
    \node (b) at (30mm,65mm) {$\frac{1}{8}$};
    \node (c) at (50mm,65mm) {$\frac{1}{16}$};
    \node (d) at (10mm,0mm) {$\frac{1}{64}$};
    \node (e) at (30mm,0mm) {$\frac{1}{4}$};
    \node (f) at (50mm,0mm) {$\frac{1}{32}$};
  \begin{scope}[->,color=gray, thick]
    \draw (a)--(19mm,34mm);
    \draw (b)--(29mm,54mm);
    \draw (c)--(40mm,35mm);
    \draw (d)--(19mm,14mm);
    \draw (e)--(29mm,35mm);
    \draw (f)--(41mm,16mm);
  \end{scope}
\end{tikzpicture}

\end{wrapfigure}

Alcune unità frazionarie più comuni venivano indicate con le parti dell'occhio 
di Horus; secondo la leggenda Horus, nella lotta contro lo zio Seth, 
reo di avergli ucciso il padre, perse un occhio le cui parti vennero ritrovate 
e ricomposte dal dio Toth a meno di una piccola parte.

I Romani fecero poco uso dei numeri frazionari; si limitarono a considerare
le parti delle misure in uso che venivano divise in~12, 24, 36, 48\ldots 
Avevano pertanto simboli e nomi particolari per indicare alcune frazioni. 
\emph{Semis} per indicare~$\frac{1}{2}$, il cui simbolo era~$S$ oppure~$Z$ 
\emph{sextans} per indicare~$\frac{1}{6}$, \emph{dracma} per 
indicare~$\frac{1}{96}$ e \emph{obolus} per indicare la sesta parte della 
\emph{dracma}.

Furono gli arabi a introdurre l'attuale scrittura delle
frazioni e i termini \emph{numeratore} e \emph{denominatore}.

La notazione attuale per le frazioni si deve sostanzialmente agli arabi, 
in Europa fu diffusa da Leonardo Pisano (Fibonacci) che con il suo 
\emph{Liber Abaci} (1202) scrive e opera con le frazioni come oggi le 
conosciamo.

\section{Frazioni}
\label{sec:03_frazioni}

\begin{definizione}
Una \emph{frazione} è una coppia ordinata di numeri naturali in cui il primo 
si chiama numeratore e il secondo denominatore. 
Il denominatore deve essere diverso da zero.
\end{definizione}

\begin{center}
 % (c) 2012 Dimitrios Vrettos - d.vrettos@gmail.com
\begin{tikzpicture}
  \begin{scope}[font=\large]
    \matrix (frazione)  at (0,0){%
      \node (num) {$a$};\\
      \node (den) {$n$};\\
    };
  \end{scope}

  \draw (num.south west)--(num.south east);

  \begin{scope}[font=\small]
    \node (testo1) at (25mm,7mm) {numeratore};
    \node (testo2) at (25mm,-7mm) {denominatore};
    \node (testo3) at (25mm,-12mm) {$n\neq0$};
  \end{scope}

  \draw[->,blue] (testo1) .. controls +(down:10mm) and +(up:10mm) .. (num) ;
  \draw[->,red] (testo2) .. controls +(up:10mm) and +(down:10mm) .. (den) ;
\end{tikzpicture}

\end{center}

Quando si chiede, per esempio un quarto di litro di 
latte,~$\frac{1}{4}\unit{l}$, si danno le informazioni su come operare 
sulla grandezza unitaria litro per ottenere la quantità desiderata.
Le frazioni possono essere viste come operatori che si applicano a una 
grandezza fissata, considerata come l'intero o il tutto, per ottenere una 
nuova grandezza ben determinata e omogenea alla prima.

Una frazione con numeratore uguale a~1 è detta \emph{frazione unitaria}; 
indicata con~$A$ una grandezza (segmento, peso, superficie, angolo\ldots) 
la scrittura~$\frac{1}{n}A$ sta ad indicare l'operazione di divisione della 
grandezza~$A$, intesa come il `tutto', in~$n$ parti uguali.

Nella figura, il segmento unitario da~0 a~1 è stato diviso in due parti uguali 
ottenendo la frazione~$\frac{1}{2}$
dividendolo in quattro parti uguali si ottiene la frazione~$\frac{1}{4}$
dividendolo in otto parti uguali si ottiene la frazione~$\frac{1}{8}$
dividendolo in sedici parti uguali si ottiene la frazione~$\frac{1}{16}$.

\begin{center}
 % (c) 2012 Dimitrios Vrettos - d.vrettos@gmail.com
% 
\begin{tikzpicture}[decoration={markings,mark=between positions 0.7 and .9 step 90pt with {\arrow{stealth}}}]
\begin{scope}[thick]
\draw[->] (0,0) -- (240pt,0);
  \foreach \x/\xtext in {0/0,.5/\frac{1}{16},1/\frac{1}{8},1.5/,2/\frac{1}{4},2.5/,3/,3.5/,4/\frac{1}{2},4.5/,5/,5.5/,6/,6.5/,7/,7.5/,8/1}
    \draw[shift={(\x,0)}] (0pt,2pt) -- (0pt,-2pt) node[below] {$\xtext$};
\end{scope}
\begin{scope}[dotted, color=gray!50!black]
     \draw[postaction={decorate}](0,5pt) arc (180:0:0.25);
     \draw[postaction={decorate}](0,5pt) arc (180:0:0.5);
     \draw[postaction={decorate}](0,5pt) arc (180:0:1);
     \draw[postaction={decorate}](0,5pt) arc (180:0:2);
\end{scope}
\end{tikzpicture}

\end{center}

\begin{definizione}
Il \emph{denominatore} di una frazione è quel numero che indica in quante 
parti uguali si è diviso l'intero.
Poiché non ha senso dividere un intero in zero parti, il denominatore deve 
essere diverso da zero.
\end{definizione}

\begin{wrapfloat}{figure}{r}{0pt}
 % (c) 2012 Dimitrios Vrettos - d.vrettos@gmail.com
% Quadrato 1/4
\begin{tikzpicture}
  \begin{scope}[scale=1,every node/.style={minimum size=5mm}]
    \draw[step=5mm, black] (0,0) grid (1,1);
    \fill[gray, opacity=.4] (0,.5) rectangle (.5,1);
    \draw[black,thick] (0,0) rectangle (1,1);
  \end{scope} 
\end{tikzpicture}

\end{wrapfloat}

Vediamo un altro esempio. Il quadrato~$Q$ della figura è stato diviso in 
quattro parti uguali e una parte è stata colorata di grigio; 
questa parte viene indicata con la frazione unitaria~$\frac{1}{4}Q$.

Una frazione~$\frac{1}{n}A$ significa l'ennesima parte di~$A$, dove~$A$ è 
il tutto che si deve dividere in~$n$ parti uguali. In altre parole,
$A$ si può ottenere moltiplicando per~$n$ la frazione~$\frac{1}{n}A$.

Partendo da~$\frac{1}{n}A$ si possono considerare i suoi multipli interi:
\[\frac{2}{n}A, \frac{3}{n}A, \ldots, \frac{n}{n}A \] che rappresentano il
doppio di un ennesimo, il triplo di un ennesimo, l'intera grandezza~$A$.

Riferendoci all'esempio del quadrato:
\begin{center}
 % (c) 2012 Dimitrios Vrettos - d.vrettos@gmail.com
% quadrati e parti di quadrati espresse in frazioni
\begin{tikzpicture}
  \begin{scope}[scale=1,every node/.style={minimum size=5mm}]
    \draw[step=5mm, black] (0,0) grid (1cm,1cm);
    \fill[gray, opacity=.4] (0,.5cm) rectangle (1cm,1cm);
    \draw[black,thick] (0,0) rectangle (1cm,1cm);

    \draw[step=5mm, draw=black, fill=gray] (3cm,0) grid (4cm,1cm);
    \fill[gray, opacity=.4] (3cm,0) rectangle (3.5cm,1cm) ;
    \fill[gray, opacity=.4] (3.5cm,0.5cm)rectangle (4cm,1cm);
    \draw[black,thick] (3cm,0) rectangle (3cm,1cm);

    \draw[step=5mm, draw=black, fill=gray] (6cm,0) grid (7cm,1cm);
    \fill[gray, opacity=.4] (6cm,0) rectangle (7cm,1cm) ;   
    \draw[black,thick] (6cm,0) rectangle (7cm,1cm);
  \end{scope} 
  \node (a) at (1.5cm,.5cm) {$\frac{2}{4}Q$};
  \node (b) at (4.5cm,.5cm) {$\frac{3}{4}Q$};
  \node (c) at (7.5cm,.5cm) {$\frac{4}{4}Q$};
\end{tikzpicture}

\end{center}

La frazione~$\frac{m}{n}A$ (si legge \emph{emme ennesimi di}~$A$) con~$m < n$ 
indica il multiplo secondo~$m$ della frazione unitaria~$\frac{1}{n}A$
essa indica la grandezza che si ottiene dividendo~$A$ in~$n$ parti uguali 
e prendendone~$m$.

\begin{definizione}
 Il \emph{numeratore} di una frazione è quel numero che esprime quante parti, 
 dell'intero suddiviso in parti secondo il denominatore, sono state prese.
\end{definizione}

Per leggere una frazione si legge prima il numeratore e poi il denominatore.
Quest'ultimo si legge come numero ordinale (terzo, quarto, quinto,\ldots) 
fino a~10 e se è maggiore di dieci si aggiunge la terminazione \emph{-esimo}.

\begin{exrig}
\begin{esempio}
 Lettura di frazioni.
 \begin{multicols}{3}
 \begin{enumeratea}
\item $\dfrac{1}{2}$ si legge \emph{un mezzo};
\item $\dfrac{1}{10}$ si legge \emph{un decimo};
\item $\dfrac{2}{3}$, si legge \emph{due terzi};
\item $\dfrac{1}{11}$ si legge \emph{un undicesimo};
\item $\dfrac{5}{7}$ si legge \emph{cinque settimi};
\item $\dfrac{1}{12}$ si legge \emph{un dodicesimo}.
\end{enumeratea}
\end{multicols}
\end{esempio}
\end{exrig}

A volte per scrivere le frazioni si utilizza la scrittura del tipo~$a/b$, 
quindi~$2/3$~$4/6$~$6/9$\ldots

\begin{definizione}
Si chiamano \emph{proprie} le frazioni che hanno il numeratore minore del 
denominatore.
Esse rappresentano sempre una grandezza minore dell'intero.
\end{definizione}

Vi sono frazioni che pur essendo formate da numeratori e denominatori 
diversi rappresentano la stessa parte dell'intero.
\begin{center}
 % (c) 2012 Dimitrios Vrettos - d.vrettos@gmail.com
% Parti uguali dell'intero
\begin{tikzpicture}
  \draw[step=3cm, black] (0,0) grid (9cm,1cm);
  \fill[gray, opacity=.4] (0,0cm) rectangle (6cm,1cm);
  \draw[black,thick] (0,0) rectangle (9cm,1cm);

  \draw[step=1.5cm, black] (0,-.5cm) grid (9cm,-1.5cm);
  \fill[gray, opacity=.4] (0,-.5cm) rectangle (6cm,-1.5cm);
  \draw[black,thick] (0,-.5cm) rectangle (9cm,-1.5cm);

  \draw[step=1cm, black] (0,-3cm) grid (9cm,-2cm);
  \fill[gray, opacity=.4] (0,-3cm) rectangle (6cm,-2cm);
  \draw[black,thick] (0,-3cm) rectangle (9cm,-2cm);
  
  \node (a) at (-.5cm,.5cm) {$\frac{2}{3}$};
  \node (b) at (-.5cm,-1cm) {$\frac{4}{6}$};
  \node (c) at (-.5cm,-2.5cm) {$\frac{6}{9}$};
\end{tikzpicture}

\end{center}
%  \ovalbox{\risolvii \ref{ese:3.1}, \ref{ese:3.2}, \ref{ese:3.3}, \ref{ese:3.4}}

\begin{definizione}
Si dicono \emph{equivalenti} due frazioni che rappresentano la stessa parte 
dell'intero.
\end{definizione}

\begin{proprieta}[Invariantiva delle frazioni]
Se si moltiplica, o si divide, numeratore e denominatore di una stessa 
frazione per uno stesso numero diverso da zero si ottiene una frazione 
equivalente alla frazione data.
\end{proprieta}

Per trovare una frazione equivalente a una frazione assegnata è sufficiente 
moltiplicare per uno stesso numero il numeratore e il denominatore della 
frazione assegnata.

\begin{exrig}
 \begin{esempio}
 Trovare due frazioni equivalenti a~$\dfrac{4}{7}$.

Moltiplicando numeratore e denominatore per~2 si ha la frazione equivalente:
\[\frac{4\cdot2}{7\cdot2}=\frac{8}{14}.\]
Moltiplicando numeratore e denominatore per~3 si ha la frazione equivalente:
\[\frac{4\cdot3}{7\cdot3}=\frac{12}{21}.\]
 \end{esempio}
\end{exrig}

\begin{definizione}
Una frazione si dice \emph{ridotta ai minimi termini} se il numeratore e il 
denominatore sono due interi primi tra loro.
\end{definizione}

Per ridurre ai minimi termini una frazione occorre dividere numeratore e 
denominatore per il loro Massimo Comune Divisore.

\begin{exrig}
 \begin{esempio}
Ridurre ai minimi termini la frazione~$\dfrac{8}{12}$.

Scompongo in fattori~8 e~12, ottengo~$8=2^3$ e~$12=3\cdot2^2$.
Calcolo il~$\mcd$ prendendo i fattori comuni con l'esponente più piccolo; 
in questo caso~$2^2$ cioè~4. Divido numeratore e denominatore per~4:
\[\frac{8}{12} = \frac{8:4}{12:4} = \frac{2}{3}.\]
 \end{esempio}
\end{exrig}

Tutte le frazioni che hanno il denominatore (numero di parti in cui va divisa 
l'unità) uguale al numeratore (numero delle parti che vanno considerate) 
rappresentano l'intero:

\[\frac{2}{2}=\frac{3}{3}=\frac{10}{10}=1.\]

\begin{wrapfloat}{figure}{r}{0pt}
 % (c) 2012 Dimitrios Vrettos - d.vrettos@gmail.com
% Interi in forma frazionaria
\begin{tikzpicture}

\filldraw[fill=gray, opacity=.4,draw=black] (0,0) rectangle (10mm,10mm);
\draw (5mm,0) -- (5mm,10mm);

\filldraw[fill=gray, opacity=.4,draw=black] (25mm,0) rectangle (40mm,10mm);
\draw (30mm,10mm) -- (30mm,0);
\draw (35mm,10mm) -- (35mm,0);

\node (a) at (17.5mm,5mm) {$\displaystyle\frac{2}{2}=1$};
\node (a) at (47.5mm,5mm) {$\displaystyle\frac{3}{3}=1$};

\end{tikzpicture}

\end{wrapfloat}

Per esempio, se divido un quadrato in due parti uguali e ne prendo due parti 
ottengo l'intero;
se divido un quadrato in tre parti uguali e ne prendo tre parti ottengo 
l'intero,\ldots

Cosa significa costruire la grandezza~$\frac{6}{2}$ del quadrato~$Q$?
Tutte le frazioni che hanno il numeratore che è multiplo del denominatore 
rappresentano un multiplo dell'intero:
\[\frac{6}{2}=3,\qquad\frac{15}{3}=5,\qquad\frac{72}{6}=12.\]

\begin{definizione}
 Si chiamano \emph{apparenti} le frazioni che hanno il numeratore multiplo 
 del denominatore;
 esse rappresentano una grandezza multipla di quella presa come intero unitario.
\end{definizione}

Le frazioni che hanno il numeratore maggiore del denominatore rappresentano 
grandezze più grandi dell'intero.
Infatti le parti da considerare (indicate dal numeratore) sono di più delle 
parti in cui è divisa l'unità (indicate dal denominatore).
\begin{center}
 % (c) 2012 Dimitrios Vrettos - d.vrettos@gmail.com
\begin{tikzpicture}

\draw (0,0) rectangle (10mm,10mm);
\draw (0,-15mm) rectangle (10mm,-5mm);
\draw (0,-30mm) -- (0mm,-20mm)--(10mm,-20mm)--(10mm,-25mm)--%
(15mm,-25mm)--(15mm,-30mm)--(0,-30mm);

% \draw[dashed](10mm,-30mm) --(10mm,-25mm);
\begin{scope}[step=5mm, black, dashed]
\draw (0,-15mm) grid (10mm,-5mm);
\draw (0,-30mm) grid (10mm,-20mm);
\end{scope}
\node[right] (a) at (17.5mm,5mm) {I $\displaystyle\frac{5}{4}$ si ottengono dividendo il quadrato in 4 parti uguali;};
\node[right] (b) at (17.5mm,-10mm){dovendone prenderne 5 l'unità non basta.};
\node[right] (b) at (17.5mm,-25mm){La grandezza ottenuta \`e formata da  $\displaystyle\frac{4}{4}$ con l'aggiunta di  $\displaystyle\frac{1}{4}$. Cioè};
\end{tikzpicture}

\end{center}
\[\frac{5}{4}=\frac{4}{4}+\frac{1}{4}.\]

\begin{definizione}
 Si chiamano \emph{improprie} le frazioni che hanno il numeratore maggiore del 
 denominatore;
 esse rappresentano una grandezza maggiore della grandezza assegnata come intero.
\end{definizione}

% \ovalbox{\risolvii \ref{ese:3.5}, \ref{ese:3.6}, \ref{ese:3.7}, \ref{ese:3.8}, \ref{ese:3.9}, \ref{ese:3.10}, \ref{ese:3.11},
% \ref{ese:3.12}, \ref{ese:3.13}, \ref{ese:3.14}, \ref{ese:3.15}, \ref{ese:3.16}}

\section{Dalle frazioni ai numeri razionali}
\label{sec:03_razionali}

Abbiamo visto che ci sono delle frazioni che, pur essendo diverse tra di loro, 
rappresentano la stessa parte dell'intero: queste frazioni vengono chiamate 
\emph{frazioni equivalenti}.
Possiamo formare dei raggruppamenti di frazioni tra loro equivalenti, 
come nella figura \ref{fig:frazequiv}.

\begin{definizione}
Ogni raggruppamento di frazioni equivalenti è definito come un 
\emph{numero razionale assoluto} ed è rappresentato da una qualunque
frazione del raggruppamento; solitamente si sceglie la frazione ridotta ai 
minimi termini.
\end{definizione}

Nel nostro esempio~$\frac{2}{3}$ è il numero razionale rappresentante del 
raggruppamento
\[\frac{2}{3}=\bigg\lbrace\frac{2}{3},\frac{4}{6},\frac{6}{9},\frac{10}{15},
\frac{14}{21},\ldots\bigg\rbrace.\]
In questo modo abbiamo dato al simbolo~$a/b$ un nuovo significato, quello di 
numero e come tale la scrittura~$a/b$ rappresenta il quoziente indicato tra i 
due numeri naturali~$a$ e~$b$. Scriveremo~$2:3=~2/3$.

\begin{inaccessibleblock}[Figura: TODO]
 \begin{figure}[t]
\begin{center}
% (c) 2012 Dimitrios Vrettos - d.vrettos@gmail.com
\begin{tikzpicture}
  \begin{scope}[font=\scriptsize]
    %%%%%%%%%%%%%%%%%%%%%%%%
    %  POSIZIONE E DISEGNO %
    %%%%%%%%%%%%%%%%%%%%%%%%    
    \begin{scope}[shape=ellipse,thick, minimum height=30mm, minimum width=40mm]
      \node [name=a, draw=orange] at (0,0) {};
      \node [name=b, draw=blue] at (45mm,0) {};
      \node [name=c, draw=red] at (90mm,0) {};
    \end{scope}
    %%%%%%%%%%%%%%%%%%%%
    %  PARTE SUPERIORE %
    %%%%%%%%%%%%%%%%%%%%
    \begin{scope}[every node/.style={below=1mm}]
      % Primo insieme
      \node  at (a.north west) {$\displaystyle\frac{4}{6}$};
      \node at (a.north) {$\displaystyle\frac{8}{12}$};
      \node  at (a.north east){$\displaystyle\frac{16}{24}$};
      % Secondo insieme
      \node  at (b.north west) {$\displaystyle\frac{7}{6}$};
      \node at (b.north) {$\displaystyle\frac{21}{18}$};
      \node  at (b.north east){$\displaystyle\frac{35}{30}$};
      % Terzo insieme
      \node  at (c.north west) {$\displaystyle\frac{4}{2}$};
      \node at (c.north) {$\displaystyle\frac{8}{4}$};
      \node  at (c.north east){$\displaystyle\frac{30}{15}$};
    \end{scope}
    %%%%%%%%%%%%%%%%%%%
    %  PARTE CENTRALE %
    %%%%%%%%%%%%%%%%%%%
    % Primo e Secondo insieme
    \foreach \x in {a.center,b.center}
      \node at (\x) {\ldots};
    % Terzo insieme
    \node[right=5mm] at (c.center) {\ldots};
    \node[left=5mm] at (c.center) {$2$};
    %%%%%%%%%%%%%%%%%%%%
    %  PARTE INFERIORE %
    %%%%%%%%%%%%%%%%%%%%
    \begin{scope}[every node/.style={above=1mm}]
      % Primo insieme
      \node[above right] at(a.south)  {$\displaystyle\frac{18}{27}$ };
      \node[above left] at(a.south)  {$\displaystyle\frac{6}{9}$};
      \node at(a.south west) {$\displaystyle\frac{2}{3}$};
      \node at(a.south east) {$\displaystyle\frac{10}{15}$};
      % Secondo insieme
      \node at(b.south)  {$\displaystyle\frac{28}{24}$};
      \node at(b.south west) {$\displaystyle\frac{14}{12}$};
      \node at(b.south east) {$\displaystyle\frac{70}{60}$};
      % Terzo insieme
      \node at(c.south)  {$\displaystyle\frac{20}{10}$};
      \node at(c.south west) {$\displaystyle\frac{6}{3}$};
      \node at(c.south east) {$\displaystyle\frac{60}{30}$};
    \end{scope}
  \end{scope}
  %%%%%%%%%%%%%%%%%%%%
  %  TITOLI          %
  %%%%%%%%%%%%%%%%%%%%
  \begin{scope}[font=\small, every node/.style={above=8mm}]
    \node (af) at (a.30) {$\displaystyle\frac{2}{3}$};
    \node (bf) at (b.30) {$\displaystyle\frac{7}{6}$};
    \node (cf) at (c.30) {2};
  \end{scope}
  %%%%%%%%%%%%%%%%%%%%
  % COLLEGAMENTI     %
  %%%%%%%%%%%%%%%%%%%%
  \begin{scope}[thick, ->]
    \draw[orange] (a.north) .. controls +(up:5mm) and +(left:5mm) .. (af) ;
    \draw[blue] (b.north) .. controls +(up:5mm) and +(left:5mm) .. (bf) ;
    \draw[red] (c.north) .. controls +(up:5mm) and +(left:5mm) .. (cf) ;
  \end{scope}
\end{tikzpicture}

\end{center}
\label{fig:frazequiv}
\caption{Esempi di frazioni equivalenti.}
\end{figure}
\end{inaccessibleblock}

\begin{definizione}
Un numero razionale assoluto preceduto dal segno è detto 
\emph{numero razionale}.
L'insieme dei numeri razionali relativi si indica con il simbolo~$\insQ$.
\end{definizione}

Il segno del numero razionale relativo è quello che si ottiene dalla regola della
divisione dei segni tra numeratore e denominatore.

\begin{exrig}
\begin{esempio}
Segno di numeri razionali.
\[\frac{-2}{-3}=+\frac{2}{3};\qquad\frac{2}{-3}=-\frac{2}{3};\qquad
\frac{-2}{3}=-\frac{2}{3}.\]
\end{esempio}
\end{exrig}

Le frazioni proprie, che hanno numeratore minore del denominatore, 
rappresentano sempre un numero compreso tra~0 e~1.

Le frazioni improprie, che hanno numeratore maggiore del denominatore, si 
possono scrivere come somma di un numero
naturale e di una frazione propria:

\begin{itemize*}
 \item il numero naturale è il risultato della divisione intera tra numeratore 
 e denominatore;
 \item il numeratore della frazione propria è il resto della divisione tra 
 numeratore e denominatore;
 \item il denominatore della frazione propria è il denominatore stesso della 
 frazione.
\end{itemize*}

Le frazioni apparenti, del tipo
$\frac{2}{2},\frac{6}{3},\frac{20}{5},\frac{12}{4},\frac{12}{3},\ldots$
corrispondono a un numero intero, rispettivamente a~1, 2, 4, 3, 4.

\begin{exrig}
 \begin{esempio}
 ~$\dfrac{11}{3}=3+\dfrac{2}{3}$.
  \begin{itemize*}
   \item $11 \div~3 =~3$ il numero naturale;
   \item $11\bmod~3 =~2$ numeratore della frazione propria;
   \item $3 =$ denominatore della frazione propria.
  \end{itemize*}
%\[\frac{11}{3}=3+\frac{2}{3}.\]
 \end{esempio}

 \begin{esempio}
 ~$\dfrac{19}{7}=2+\dfrac{5}{7}$.
  \begin{itemize*}
   \item $19 \div~7 =~2$ il numero naturale;
   \item $19\bmod~7 =~2$ numeratore della frazione propria;
   \item $5 =$ denominatore della frazione propria.
  \end{itemize*}
%\[\frac{19}{7}=2+\frac{5}{7}.\]
 \end{esempio}
\end{exrig}

% \ovalbox{\risolvii \ref{ese:3.17}, \ref{ese:3.18}}

\section{La scrittura dei numeri razionali}
\label{sec:03_decimali}

I numeri razionali, rappresentati finora come frazioni, possono essere scritti 
come numeri decimali:
basta fare la divisione tra numeratore e denominatore, il quoziente ottenuto 
è la rappresentazione della frazione sotto forma decimale.

\begin{center}
 % (c) 2012 Dimitrios Vrettos - d.vrettos@gmail.com
\begin{tikzpicture}
  \begin{scope}[font=\ttfamily\small,  ]
    \matrix (a) [matrix of nodes]{
    1 &{}&{}&{}&{}&3&{}&{}&{}&{}&{}&{}&[20mm]&1&1&{}&{}&{}&8\\
    1 &0 &{}&{}&{}&|[blue]|0&|[blue]|,&|[blue]|3&|[blue]|3&|[blue]|3&|[blue]|3&|[font=\scriptsize,blue]|\ldots&{}&{}&3&{}&{}&{}&|[blue]|1&|[blue]|,&|[blue]|3&|[blue]|7&|[blue]|5\\
    {}&1 &0 &{}&{}&{}&{}&{}&{}&{}&{}&{}&{}&{}&3&0\\
    &&1&0&&&{}&{}&{}&{}&{}&{}&{}&{}&{}&6&0\\
    &&&1&0&&{}&{}&{}&{}&{}&{}&{}&{}&{}&{}&4&0\\
    &&&&|[font=\scriptsize, red]|\ldots&{}&&{}&{}&{}&{}&{}&{}&{}&{}&{}&{}&|[red]|0\\
    };
  \end{scope}
  \begin{scope}[densely dotted,->]
    \draw (a-1-2.center) -- (a-2-2);
    \draw (a-1-3.center) -- (a-3-3);
    \draw (a-1-4.center) -- (a-4-4);
    \draw (a-1-5.center) -- (a-5-5);

    \draw (a-1-16.center) -- (a-3-16);
    \draw (a-1-17.center) -- (a-4-17);
    \draw (a-1-18.center) -- (a-5-18);
  \end{scope}

  \draw (a-1-6.north west) -- (a-1-6.south west);
  \draw (a-2-6.north west) -- (a-2-11.north east);

  \draw (a-1-19.north west) -- (a-1-19.south west);
  \draw (a-2-19.north west) -- (a-2-23.north east);

  \begin{scope}[font=\small]
    \node[below] at(a-6-6.south) {$\displaystyle\frac{1}{3}=0,3333\ldots$};
    \node[below=3mm] at(a-6-18.east) {$\displaystyle\frac{11}{8}=1,375$};
  \end{scope}
\end{tikzpicture}

\end{center}

I numeri decimali che si ottengono sono di due tipi: numeri decimali finiti 
come~$1,375$ e numeri decimali periodici come~$1,333333\ldots$
quest'ultimo si scrive mettendo una barra sulla parte 
periodica:~$1,\overline{3}$ oppure racchiudendo la parte periodica tra parentesi 
tonde~$1,(3)$.

I numeri decimali finiti si ottengono dalle frazioni il cui denominatore ha 
come fattori solo il~2, solo il~5 o entrambi, eventualmente elevati a una 
potenza.

I numeri decimali periodici semplici si ottengono dalle frazioni il cui 
denominatore non ha per fattori né~2 né~5.

I numeri decimali periodici misti si ottengono dalle frazioni il cui 
denominatore contiene altri fattori oltre al~2 e al~5.

\begin{exrig}
 \begin{esempio}
 Alcuni numeri decimali finiti.
\begin{enumeratea}
 \item $\dfrac{11}{8}=\dfrac{11}{2^3}=
\dfrac{11\cdot 5^3}{2^3\cdot 5^3}=\dfrac{1375}{1000}=1,375$
 \item $\dfrac{7}{25}=\dfrac{7}{5^2}=
\dfrac{7\cdot 2^2}{5^2\cdot 2^2}=\dfrac{28}{100}=0,28~$
 \item $\dfrac{13}{40}=\dfrac{13}{2^3\cdot 5}=
\dfrac{13\cdot 5^2}{2^3\cdot 5^3}=\dfrac{325}{1000}=0,325$
 \item $\dfrac{50}{7} = \dfrac{\ldots}{10}$, non è possibile, non è un 
 decimale finito.
\end{enumeratea}
 \end{esempio}
\end{exrig}

% \ovalbox{\risolvi \ref{ese:3.19}}

\begin{procedura}
Trasformare una frazione in numero decimale:
\begin{enumeratea}
 \item eseguire la divisione tra numeratore e denominatore;
 \item se la divisione ha un resto mettere la virgola al quoziente e 
 moltiplicare per~$10$ il resto;
 \item continuare la divisione finché il resto è zero oppure fino a che non 
 si trova un resto già trovato prima;
 \item se la divisione si conclude con resto~$0$ si ottiene un numero decimale 
 finito;
 \item se la divisione si conclude perché si è ritrovato un resto ottenuto in 
 precedenza si ottiene un numero decimale periodico.
\end{enumeratea}
\end{procedura}

\clearpage

\begin{exrig}
 \begin{esempio}
 Trasformazione di frazioni in numeri decimali.
  \begin{center}
  % (c) 2012 Dimitrios Vrettos - d.vrettos@gmail.com
% Da frazione al decimale

\begin{tikzpicture}
  \begin{scope}[font=\ttfamily\scriptsize,matrix anchor=north west]

    \matrix at (0,0) [name= a, matrix of nodes]{
    {}&1&1&3&{}&{}&2&0\\
    -&1&0&0&{}&{}&|[blue]|5&|[blue]|,&|[blue]|6&|[blue]|5\\
    {}&{}&1&3&0\\
    {}&-&1&2&0\\
    {}&{}&{}&1&0&0\\
    {}&{}&-&1&0&0\\
    {}&{}&{}&{}&{}&|[red]|0\\};

    \matrix at(39mm,0) [name=b, matrix of nodes]{
    &{}&1&7&{}&{}&6\\
    &-&1&2&{}&{}&|[blue]|2&|[blue]|,&|[blue]|8&|[blue]|3\\
    &{}&{}&5&0\\
    &{}&-&4&8\\
    &{}&{}&{}&2&0\\
    &{}&{}&-&1&8\\
    &{}&{}&{}&{}&|[red]|2\\
    };

    \matrix at (76mm,0) [name=c, matrix of nodes] {
    &{}&1&5&{}&{}&{}&{}&{}&{}&7\\
    &-&1&4&{}&{}&{}&{}&{}&{}&|[blue]|2&|[blue]|,&|[blue]|1&|[blue]|4&|[blue]|2&|[blue]|8&|[blue]|5&|[blue]|7\\
    &{}&{}&1&0\\
    &{}&-&{}&7\\
    &{}&{}&{}&3&0\\
    &{}&{}&-&2&8\\
    &{}&{}&{}&{}&2&0\\
    &{}&{}&{}&-&1&4\\
    &{}&{}&{}&{}&{}&6&0\\
    &{}&{}&{}&{}&-&5&6\\
    &{}&{}&{}&{}&{}&{}&4&0\\
    &{}&{}&{}&{}&{}&-&3&5\\
    &{}&{}&{}&{}&{}&{}&{}&5&0\\
    &{}&{}&{}&{}&{}&{}&-&4&9\\
    &{}&{}&{}&{}&{}&{}&{}&{}&|[red]|1\\
    };

  \end{scope}
  % Prima divisione
  \draw (a-1-7.north west) -- (a-2-7.north west);
  \draw (a-2-7.north west) -- (a-2-10.north east);
  \draw (a-2-2.south west) -- (a-2-4.south east);
  \draw (a-4-3.south west) -- (a-4-5.south east);
  \draw (a-6-4.south west) -- (a-6-6.south east);

  % Seconda divisione
  \draw (b-1-7.north west) -- (b-2-7.north west);
  \draw (b-2-7.north west) -- (b-2-10.north east);
  \draw (b-2-3.south west) -- (b-2-4.south east);
  \draw (b-4-4.south west) -- (b-4-5.south east);
  \draw (b-6-5.south west) -- (b-6-6.south east);

  % Terza divisione
  \draw (c-1-11.north west) -- (c-2-11.north west);
  \draw (c-2-11.north west) -- (c-2-18.north east);
  \draw (c-2-3.south west) -- (c-2-4.south east);
  \draw (c-4-4.south west) -- (c-4-5.south east);
  \draw (c-6-5.south west) -- (c-6-6.south east);
  \draw (c-8-6.south west) -- (c-8-7.south east);
  \draw (c-10-7.south west) -- (c-10-8.south east);
  \draw (c-12-8.south west) -- (c-12-9.south east);
  \draw (c-14-9.south west) -- (c-14-10.south east);

  \begin{scope}[dashed, densely dotted, ->]
    \draw (a-1-5.center) -- (a-3-5);
    \draw (a-1-6.center) -- (a-5-6);

    \draw (b-1-5.center) -- (b-3-5);
    \draw (b-1-6.center) -- (b-5-6);

    \draw (c-1-5.center) -- (c-3-5);
    \draw (c-1-6.center) -- (c-5-6);
    \draw (c-1-7.center) -- (c-7-7);
    \draw (c-1-8.center) -- (c-9-8);
    \draw (c-1-9.center) -- (c-11-9);
    \draw (c-1-10.center) -- (c-13-10);
  \end{scope}
  \begin{scope}[font=\small]
    \node[above=5pt]  at (a-1-6.north) {a)};
    \node[above=5pt]  at (b-1-6.north) {b)};
    \node[above=5pt]  at (c-1-10.north) {c)};
  \end{scope}
\end{tikzpicture}


  \end{center}
  \begin{enumeratea}
  \item $\dfrac{113}{20}=5,65$, numero decimale finito;\vspace{1.03ex}
  \item $\dfrac{17}{6}=2,8\overline{3}$, numero decimale periodico misto di 
  periodo~3;\vspace{1.03ex}
  \item $\dfrac{15}{7}=2,\overline{142857}$, numero decimale periodico di 
  periodo~142857.
  \end{enumeratea}
 \end{esempio}
\end{exrig}


% \ovalbox{\risolvi \ref{ese:3.20}, \ref{ese:3.21}}\vspazio

Viceversa un numero decimale finito o periodico può essere sempre scritto 
sotto forma di frazione.

\begin{procedura}
Trasformare un numero decimale finito in una frazione:
\begin{enumeratea}
 \item scrivere una frazione che ha per numeratore il numero che si vuole 
  trasformare e per denominatore~uno;
 \item moltiplicare numeratore e denominatore per dieci elevato al numero
  di cifre a destra della virgola.
 \item semplificare la frazione così ottenuta.
\end{enumeratea}
\end{procedura}

Per facilitare questa operazione possiamo considerare i numeri decimali finiti 
come frazioni particolari che hanno il numeratore uguale al numero decimale e 
il denominatore uguale a~1.

\begin{exrig}
\begin{esempio}
Trovare la frazione equivalente a~1,36:
\[\frac{1,36}{1}=\frac{1,36\cdot10^2}{1\cdot10^2}=\frac{136}{100}=\frac{34}{25}.\].
\end{esempio}
\begin{esempio}
Trovare la frazione equivalente a~0,00043000:
\[\frac{0,00043}{1}=\frac{0,00043\cdot10^5}{1\cdot10^5}=\frac{43}{100000}.\]
\end{esempio}
\end{exrig}

Un numero decimale periodico, generalmente, presenta tre elementi:
\begin{description}
 \item [la parte intera] composta dalle cifre poste prima della virgola;
 \item [il periodo] che è composto da una o più cifre che si ripetono 
 all'infinito dopo la virgola;
 \item [l'antiperiodo] la parte composta da zero o più cifre  poste tra la 
 virgola e il periodo.
\end{description}


Per esempio, nel numero~253,485795795795795\ldots la parte intera è~253, 
il periodo è~579, l'antiperiodo è~48.

Dato che il numero è infinito non può essere scritto con tutte le sue cifre, 
si usano due modi per scriverlo in forma compatta, mettendo una lineetta 
sopra le cifre del periodo o racchiudendo le cifre del periodo tra parentesi 
tonde.

Il numero~253,485795795795795\ldots può essere scritto~$253,48\overline{579}$, 
oppure~$253,48(579)$.

I numeri decimali periodici si dividono in:
\begin{description}
 \item [semplici] se subito dopo la virgola è presente il periodo;
 \item [misti] se dopo la virgola è presente l'antiperiodo.
\end{description}

Anche i numeri periodici possono essere trasformati in una frazione, che si
dice \emph{frazione generatrice} del numero.

\begin{procedura}
 Determinare la frazione generatrice di un numero periodico:
\begin{enumeratea}
 \item il numeratore della frazione si ottiene sottraendo dal numero 
  senza la virgola e con il periodo scritto una sola volta,
  il numero costituito dalle cifre che precedono il periodo;
 \item il denominatore della frazione si ottiene scrivendo tanti~$9$ 
  quante sono le cifre del periodo seguiti da
  tanti~$0$ quante sono le cifre comprese tra il periodo e la virgola;
 \item semplificare la frazione ottenuta.
\end{enumeratea}
\end{procedura}

\paragraph*{Passo~$a$}~$2,5\overline{12}\rightarrow~2512$.
\paragraph*{Passo~$b$}~$2512-25=2487$.
\paragraph*{Passo~$c$}~$2,5\overline{12}=\dfrac{2487}{990}$.


\paragraph*{Ma perché questa regola? Una possibile spiegazione}
Consideriamo il numero periodico semplice~$2,\overline{3}$. Considero la 
frazione~$\frac{2,\overline{3}}{1}$
moltiplico numeratore e denominatore 
per~10~$\frac{2,\overline{3}\cdot10}{1\cdot10}$ e ottengo
$\frac{23,\overline{3}}{10}$.

L'obiettivo è quello di eliminare dal numeratore della frazione la parte 
decimale.
Per ottenere questo risultato tolgo~$2,\overline{3}$ da~$23,\overline{3}$, 
cioè $23,\overline{3}-2,\overline{3}=21$.

Come mai~$2,\overline{3}$ e non~$1,\overline{3}$ o~$0,\overline{3}$?
Perché in questo modo posso sapere quanto vale il denominatore: 
se~$23,\overline{3}$ è il risultato
della moltiplicazione di~$2,\overline{3}\cdot10$,~$21$~è il risultato 
della moltiplicazione di
$2,\overline{3}\cdot9$ in quanto~$23,\overline{3}-2,\overline{3}=21$.
In definitiva
\[2,\overline{3}=\frac{23-2}{9}=\frac{21}{9}=\frac{7}{3}.\]

Possiamo usare lo stesso procedimento per il numero periodico 
misto~$2,5\overline{12}$.

Considero la frazione~$\frac{2,5\overline{12}}{1}$, moltiplico numeratore 
e denominatore per~1000 e ottengo: $\frac{2512,\overline{12}}{1000}$.
L'obiettivo è quello di eliminare dal numeratore della frazione la parte 
decimale che contiene il periodo che si ripete all'infinito. 
Per ottenere questo risultato tolgo da~$2512,\overline{12}$ questa volta
$25,\overline{12}$, cioè~$2512,\overline{12}-25,\overline{12}=2487$.
Per avere una frazione equivalente occorre che al denominatore abbia~990 in 
quanto dal numeratore ho tolto~10 volte $2,5\overline{12}$.
\[2,5\overline{12}=\frac{2512-25}{990}=\frac{2487}{990}.\]

% \ovalbox{\risolvii \ref{ese:3.22}, \ref{ese:3.23}}

\subsection{Numeri periodici particolari}

Numeri periodici particolari sono quelli che hanno come periodo il numero~9,
come~$2,\overline{9}$, $1,1\overline{9}$, $21,22\overline{9}$ ecc.
Se, per esempio, applichiamo la regola per il calcolo della frazione 
generatrice al numero periodico otteniamo un risultato inatteso
\[2,\overline{9}=\frac{29-2}{9}=\frac{27}{9}=3.\]
Quindi~$2,\overline{9}$ coincide con il numero intero~3.

Per lo stesso motivo~$1,1\overline{9}=2$,~$21,22\overline{9}=21,23$.

\begin{wrapfloat}{figure}{r}{0pt}
% (c) 2012 Dimitrios Vrettos - d.vrettos@gmail.com
\begin{tikzpicture}
  \begin{scope}[thick,font=\scriptsize]
  \draw[->] (0,0) -- (20mm,0);
  \foreach \n in {5mm,15mm}{%
    \draw (\n,-3pt) -- (\n,3pt);}
  \node (p) at (5mm,-4mm) {$0,\overline{9}$};
  \node (i) at (15mm,-4mm) {$1$};
  \end{scope}
\end{tikzpicture}

\end{wrapfloat}

Questo fatto si può anche dimostrare in modo grafico, rappresentando, 
ad esempio, il numero~$0,\overline{9}$ e il numero~1 sulla retta reale.

Se i due numeri fossero veramente diversi sarebbero rappresentati da due punti 
distinti come in figura. Dato che la retta reale non può avere ``buchi'',
tra un suo punto e un altro ci deve essere almeno un altro numero compreso tra 
i due.
Ma qual è questo numero? Qualunque numero decimale minore di~1 è sicuramente 
superato dal numero~$0,\overline{9}$,
ad esempio~0,9999999998 è sicuramente più piccolo di~$0,\overline{9}$. 
Quindi non esiste nessun numero tra $0,\overline{9}$ e~1,
di conseguenza i due numeri coincidono.

% \ovalbox{\risolvii \ref{ese:3.24}, \ref{ese:3.25}}

\section{I numeri razionali e la retta}
\label{sec:03_retta}

Anche i numeri razionali si possono rappresentare su una retta orientata. 
Per fare questo occorre scegliere un punto~$O$ sulla retta e associare ad esso 
il numero zero. Fissiamo poi un segmento unitario e scegliamo un verso di 
percorrenza.

Dato un numero razionale positivo, rappresentato dalla frazione~$\frac{a}{n}$, 
il punto corrispondente al numero razionale sulla retta viene determinato nel 
seguente modo. Dividiamo il segmento unitario~$u$ in tante parti uguali
quante sono quelle indicate dal denominatore~$n$ della frazione, ottenendo 
così la frazione unitaria~$\frac{1}{n}$.
A partire dal punto~$O$ procedendo verso destra, si contano~$a$ frazioni 
unitarie. L'ultimo punto rappresenta il numero razionale~$\frac{a}{n}$.

Per le frazioni improprie la singola unità~$u$ non è sufficiente, occorre 
prendere la unità successiva di~$u$ e dividere anche questa in~$n$ parti. 
Il procedimento si ripete fino a che si considerano tutte le frazioni unitarie 
indicate da~$a$. Anche in questo caso, il punto individuato dall'ultima 
frazione unitaria rappresenta il numero razionale~$\frac{a}{n}$. 
In alternativa si può scomporre la frazione impropria nella somma di un numero 
intero e di una frazione propria, quindi si rappresenta la frazione impropria
a partire dal suo numero intero invece che partire da~0. Per esempio, per 
rappresentare la frazione~$\frac{3}{2}$ trasformiamo la frazione 
in~$1+\frac{1}{2}$, quindi rappresentiamo partendo dal numero~1 invece che 
da~0.

Se il numero razionale è negativo, ci comportiamo come prima con l'avvertenza 
di muoverci nel senso opposto a quello precedente cioè da destra verso 
sinistra.

\begin{center}
% (c) 2012 Dimitrios Vrettos - d.vrettos@gmail.com
\begin{tikzpicture}
\begin{scope}[thick,font=\small]
\draw[->] (-10pt,0) -- (240pt,0) node [above] () {$\insQ$};
  \foreach \x/\xtext in {%
      0/-2,.25/,.5/,.75/-\frac{13}{8},1/,1.25/,1.5/,1.75/,2/-1,2.25/,2.5/,2.75/,3/-\frac{1}{2},%
      3.25/,3.5/,3.75/,4/0,4.25/,4.5/,4.75/\frac {3}{8},%
      5/,5.25/,5.5/,5.75/,6/1,6.25/,6.5/,6.75/,7/\frac{3}{2},7.25/,7.5/,7.75/,8/2}
    \draw[shift={(\x,0)}] (0pt,2pt) -- (0pt,-2pt) node[below] {$\xtext$};
\draw (.75,0) circle (1pt);
\draw (3,0) circle (1pt);
\draw (4.75,0) circle (1pt);
\draw (7,0) circle (1pt);
\end{scope}
\end{tikzpicture}

\end{center}

% \ovalbox{\risolvii \ref{ese:3.26}, \ref{ese:3.27}, \ref{ese:3.28}}

\section{Confronto tra numeri razionali}
\label{sec:03_confronto}

Il numero razionale rappresentato dalla frazione~$\frac{a}{n}$ è \emph{minore} 
del numero razionale rappresentato dalla frazione~$\frac{b}{m}$, se nella 
retta orientata il punto che corrisponde alla frazione$\frac{a}{n}$ precede il 
punto che corrisponde alla frazione$\frac{b}{m}$ e si scrive
\[\frac{a}{n}<\frac{b}{m}.\]

Il numero razionale~$\frac{a}{n}$ è \emph{maggiore} di~$\frac{b}{m}$,
se nella retta orientata il punto che corrisponde alla frazione~$\frac{a}{n}$ 
segue il punto che corrisponde alla frazione~$\frac{b}{m}$ e si scrive
\[\frac{a}{n}>\frac{b}{m}.\]

Il numero razionale~$\frac{a}{n}$ è \emph{equivalente} a~$\frac{b}{m}$ se 
nella retta orientata i punti che corrispondono
alle frazioni~$\frac{a}{n}$ e~$\frac{b}{m}$ coincidono.

%\newpage
\begin{exrig}
\begin{esempio}
Confronto tra numeri razionali.
\begin{center}
% (c) 2012 Dimitrios Vrettos - d.vrettos@gmail.com
\begin{tikzpicture}
\begin{scope}[thick,font=\small]
\draw[->] (-10pt,0) -- (240pt,0) node [above] () {$\insQ$};
  \foreach \x/\xtext in {%
      0/-2,.25/,.5/,.75/-\frac{13}{8},1/,1.25/,1.5/,1.75/,2/-1,2.25/,2.5/,2.75/,3/-\frac{1}{2},%
      3.25/,3.5/,3.75/,4/0,4.25/,4.5/,4.75/\frac {3}{8},%
      5/,5.25/,5.5/,5.75/,6/1,6.25/,6.5/,6.75/,7/\frac{3}{2},7.25/,7.5/,7.75/,8/2}
    \draw[shift={(\x,0)}] (0pt,2pt) -- (0pt,-2pt) node[below] {$\xtext$};
\draw (.75,0) circle (1pt);
\draw (3,0) circle (1pt);
\draw (4.75,0) circle (1pt);
\draw (7,0) circle (1pt);
\end{scope}
\end{tikzpicture}

\end{center}
\[-\frac{13}{8}<-\frac{1}{2},\qquad\frac{3}{8}>-\frac{1}{2},\qquad
\frac{3}{8}<\frac{3}{2},\qquad-1>-\frac{13}{8}.\]
\end{esempio}
\end{exrig}

Per certe frazioni è facile vedere se una frazione precede o segue un'altra. 
Per altre non è così semplice.

Consideriamo per esempio le frazioni~$\frac{7}{9}$ e~$\frac{6}{7}$.
Quale frazione precede e quale segue? Il confronto non è immediato perché 
con la prima frazione si conta per unità frazionarie di tipo~$\frac{1}{9}$, 
con la seconda per unità frazionarie di tipo~$\frac{1}{7}$.

In generale, senza ricorrere alla rappresentazione sulla retta, come si 
possono confrontare i numeri razionali?

Conviene sostituire le frazioni date con altre equivalenti che hanno unità 
frazionarie dello stesso tipo:
cioè occorre ridurre le frazioni allo stesso denominatore.

\begin{procedura}
Confrontare due frazioni:
\begin{enumeratea}
\item si calcola il minimo comune multiplo dei denominatori delle frazioni;
\item si trasforma ciascuna frazione come segue:
\begin{itemize*}
\item il nuovo denominatore è il~$\mcm$ trovato;
\item il nuovo numeratore si ottiene dividendo il~$\mcm$ per il 
denominatore della frazione data e moltiplicando il quoziente ottenuto per il 
numeratore della frazione data.
\end{itemize*}
\item si confrontano i nuovi numeratori: la frazione più grande è quella che 
ha il numeratore più grande.
\end{enumeratea}
\end{procedura}

Un altro modo per confrontare due frazioni consiste nel 
\emph{moltiplicare in croce} numeratori e denominatori delle frazioni, 
come nei seguenti esempi.

\begin{exrig}
\begin{esempio}
Confronta~$\frac{3}{2}$ con~$\frac{5}{3}$.

Moltiplichiamo il numeratore della prima frazione con il denominatore della 
seconda frazione e il denominatore della prima frazione per il denominatore 
della seconda, così:
 \[\frac{3}{2}<\frac{5}{3}\text{, perché }3\cdot3<2\cdot5. \]
\end{esempio}

\begin{esempio}
Confronta le frazioni~$\frac{7}{9}$ e~$\frac{6}{7}$.

$\mcm(7.9)=63$.
\[\frac{7}{9}=\frac{7\cdot7}{9\cdot7}=\frac{49}{63},\qquad%
\frac{6}{7}=\frac{6\cdot9}{7\cdot9}=\frac{54}{63}.\]
\[\frac{54}{63}>\frac{49}{63}\Rightarrow\frac{6}{7}>\frac{7}{9}.\]
\end{esempio}
\end{exrig}

% \ovalbox{\risolvii \ref{ese:3.29}, \ref{ese:3.30}, \ref{ese:3.31}, \ref{ese:3.32}, \ref{ese:3.33}, \ref{ese:3.34}, \ref{ese:3.35}, \ref{ese:3.36}, \ref{ese:3.37}, \ref{ese:3.38},%
% \ref{ese:3.39}, \ref{ese:3.40}, \ref{ese:3.41},}

% \vspazio\ovalbox{\ref{ese:3.42}, \ref{ese:3.43}}

\section{Le operazioni con i numeri razionali}
\label{sec:03_operazioni}

Con i numeri razionali è sempre possibile eseguire le addizioni, 
le moltiplicazioni, le sottrazioni e le divisioni.
In altre parole, poiché un numero razionale può essere scritto sotto forma 
di frazione, se si addizionano,
si moltiplicano, si sottraggono, si dividono due frazioni il risultato è 
sempre una frazione.

\subsection{Addizione}

Se due frazioni hanno la stessa unità frazionaria allora è sufficiente sommare 
i numeratori delle frazioni e prendere come denominatore l'unità frazionaria 
comune.
\[\frac{5}{3}+\frac{2}{3}=\frac{5+2}{3}=\frac{7}{3}.\]
%
%\begin{center}
%% (c) 2012 Dimitrios Vrettos - d.vrettos@gmail.com
\begin{tikzpicture}
  \draw (0,0) rectangle  (10mm,10mm);
  \begin{scope}[font=\Huge]
    \node at  (5mm,5mm) {$+$};
  \end{scope}
  \begin{scope}[font=\small]
    \node (a) at (-15mm,10mm) {$\displaystyle\frac{5}{3}$};
    \node (b) at (-15mm,0mm) {$\displaystyle\frac{2}{3}$};
    \node (c) at (25mm,5mm) {$\displaystyle\frac{7}{3}$};
  \end{scope}
  \node (entrata) at (0,5mm) {};
  \node (uscita) at (10mm,5mm) {};
  \begin{scope}[->,concept connection,blue,thin]
    \draw (a) -- (entrata);
    \draw (b) -- (entrata);
    \draw (uscita) -- (c);
  \end{scope}
\end{tikzpicture}

%\end{center}

\begin{definizione}
La \emph{somma di due frazioni con lo stesso denominatore} è una frazione che 
ha per denominatore lo stesso denominatore delle frazioni date e per 
numeratore la somma dei numeratori.
\end{definizione}


Se le unità frazionarie sono diverse dobbiamo considerare frazioni equivalenti 
a quelle date che abbiano la stessa unità frazionaria e poi eseguire 
l'addizione come indicato nel punto precedente e cioè sommando i numeratori 
e lasciando lo stesso denominatore comune.

\begin{center}
% (c) 2012 Dimitrios Vrettos - d.vrettos@gmail.com
\begin{tikzpicture}
  \draw (0,0) rectangle  (10mm,10mm);
  \begin{scope}[font=\Huge]
    \node at  (5mm,5mm) {$+$};
  \end{scope}
  \begin{scope}[font=\small]
    \node [left](a) at (-12.5mm,10mm) {$\displaystyle{\frac{5}{3}=\frac{25}{15}}$};
    \node [left](b) at (-12.5mm,0mm) {$\displaystyle{\frac{2}{3}=\frac{6}{15}}$};
    \node (c) at (25mm,5mm) {$\displaystyle\frac{31}{15}$};
  \end{scope}
  \node (entrata) at (0,5mm) {};
  \node (uscita) at (10mm,5mm) {};
  \begin{scope}[->,concept connection,blue,thin]
    \draw (a) -- (entrata);
    \draw (b) -- (entrata);
    \draw (uscita) -- (c);
  \end{scope}
\end{tikzpicture}

\end{center}

In generale data l'addizione di due frazioni~$\frac{m}{n}+\frac{p}{q}$ la 
somma si può scrivere come
\[\frac{mq +pn}{nq}.\]

\begin{center}
% (c) 2012 Dimitrios Vrettos - d.vrettos@gmail.com
\begin{tikzpicture}
  \draw (0,0) rectangle  (10mm,10mm);
  \begin{scope}[font=\Huge]
    \node at  (5mm,5mm) {$+$};
  \end{scope}
  \begin{scope}[font=\small]
    \node [left](a) at (-12.5mm,10mm) {$\displaystyle{\frac{m}{n}=\frac{mq}{nq}}$};
    \node [left](b) at (-12.5mm,0mm) {$\displaystyle{\frac{p}{q}=\frac{pn}{nq}}$};
    \node [right](c) at (25mm,5mm) {$\displaystyle\frac{mq+pn}{nq}$};
  \end{scope}
  \node (entrata) at (0,5mm) {};
  \node (uscita) at (10mm,5mm) {};
  \begin{scope}[->,concept connection,blue,thin]
    \draw (a) -- (entrata);
    \draw (b) -- (entrata);
    \draw (uscita) -- (c);
  \end{scope}
\end{tikzpicture}

\end{center}

Quando si sommano due frazioni si può scegliere un qualsiasi denominatore 
comune, tuttavia per semplificare i calcoli conviene scegliere il più piccolo 
possibile, cioè il minimo comune multiplo dei denominatori delle frazioni da 
sommare.

\begin{procedura}
Sommare due o più frazioni:
\begin{enumeratea}
\item ridurre le frazioni ai minimi termini;
\item calcolare il~$\mcm$ dei denominatori;
\item mettere il~$\mcm$ come denominatore della frazione somma;
\item per ogni frazione dividere il~$\mcm$ per il suo denominatore e 
moltiplicare il risultato per il numeratore della frazione mantenendo il segno;
\item calcolare la somma algebrica di tutti i numeri trovati;
\item mettere la somma ottenuta come numeratore della frazione somma;
\item ridurre ai minimi termini la frazione ottenuta.
\end{enumeratea}
\end{procedura}

\begin{exrig}
\begin{esempio}
Sommare le frazioni~$\frac{8}{12}- \frac{5}{6}+\frac{8}{5}-1$.

\paragraph*{Passo~$a$} riduco ai minimi termini le frazioni
$\displaystyle{\frac{2}{3}-\frac{5}{6}+\frac{8}{5}-\frac{1}{1}}$
\paragraph*{Passo~$b$} calcolo~$\mcm(3,6,5,1)=30$.
\paragraph*{Passo~$c$} la frazione somma avrà come denominatore il~$\mcm$ 
trovato~$\dfrac{\ldots}{30}$.
\paragraph*{Passo~$d$} per ogni frazione divido il~$\mcm$ per il suo 
denominatore e moltiplico il risultato per il numeratore:
\begin{align*}
\frac{2\cdot(30:3)-5\cdot(30:6)+8\cdot(30:5)-1\cdot(30:1)}{30}&=
\frac{2\cdot10-5\cdot5+8\cdot6-1\cdot30}{30}\\
&=\frac{20-25+48-30}{30}.
\end{align*}

\paragraph*{Passo~$e$} calcolo la somma algebrica dei numeri ottenuti al 
numeratore~$+13$.
\paragraph*{Passo~$f$} metto la somma ottenuta al numeratore della frazione 
somma~$+\dfrac{13}{30}$.
\paragraph*{Passo~$g$} vedo se posso ridurre la frazione, in questo caso no, 
il risultato è~$+\dfrac{13}{30}$.
\end{esempio}

\begin{esempio}
Sommare i numeri razionali~$-0,2-1,\overline{2}+25\%+\frac{7}{12}$.

Trasformo i numeri razionali in frazioni:
\[-\frac{2}{10}-\frac{12-1}{9}+\frac{25}{100}+\frac{7}{12}=-\frac{1}{5}-
\frac{11}{9}+\frac{1}{4}+\frac{7}{12}.\]

Quindi~$\mcm(5,9,4,12)=180$.

\begin{align*}
\frac{-1\cdot(180:5)-11\cdot(180:9)+1\cdot(180:4)+7\cdot(180:12)}{180}&=
\frac{-1\cdot36-11\cdot20+1\cdot45+7\cdot15}{180}\\
&=\frac{-36-220+45+105}{180}\\
&=-\frac{106}{180} \\
&=-\frac{53}{90}.
\end{align*}
\end{esempio}
\end{exrig}

\subsection{Sottrazione di frazioni}

La sottrazione di frazioni si può sempre trasformare in una addizione tra la 
prima frazione e l'opposto della seconda frazione. 
Come per i numeri relativi, quando si parla di somma di frazioni si intende 
sempre somma algebrica di frazioni.

% \vspazio\ovalbox{\risolvii \ref{ese:3.44}, \ref{ese:3.45}, \ref{ese:3.46}, \ref{ese:3.47}}

\subsection{Moltiplicazione}

Il risultato della moltiplicazione tra frazioni può essere interpretato come 
l'area di un rettangolo in cui le frazioni fattori sono la base e l'altezza.
\begin{center}
 % (c) 2012 Dimitrios Vrettos - d.vrettos@gmail.com
\tikzset{brace/.style={decoration={brace, mirror},decorate}}

\begin{tikzpicture}[font=\small]
  \begin{scope}[x=6mm,y=10mm]
    \foreach \x in {0,1,2,3,4}
      \foreach \y in {0,1,2}
	\draw (\x,\y) rectangle (\x+1,\y+1);
    \foreach \x in {0,1,...,3}
      \foreach \y in {0,1}{
	\fill[gray, opacity=.4](\x,\y) rectangle (\x+1,\y+1);
	\draw (\x+.5,\y+.5) node [rectangle] () {$\frac{1}{5}$};
      }

    \draw (4.5,2.5) node [rectangle] () {$\frac{1}{5}$};

    \begin{scope}[|-|,above]
      \draw (0,3.2) -- (5,3.2) node [pos=.5] () {$1\text{ unità}$};
      \draw (5.33,0)--(5.33,3)node [rotate=-90,pos=.5] () {$1\text{ unità}$};
    \end{scope}

    \draw[brace] (0,-.1) --  (4,-.1) node [below, pos=0.5] {$\frac{4}{5}$};
    \draw[brace] (-.2,2) --  (-.2,0) node [left, pos=0.5] {$\frac{2}{3}$};
  \end{scope}
  
\end{tikzpicture}

\end{center}

Moltiplicare~$\frac{4}{5}\cdot\frac{2}{3}$ è come calcolare l'area del 
rettangolo di base~$\frac{4}{5}$ e altezza $\frac{2}{3}$. 
Ogni rettangolino di base~$\frac{1}{5}$ e altezza~$\frac{1}{3}$ ha 
area~$\frac{1}{15}$.
I rettangolini da prendere in considerazione sono~8. Il risultato è 
quindi~$\frac{8}{15}$. 
Il denominatore indica in quante parti è stato diviso il quadrato unitario: 
sono $3 \cdot 5=15$ parti.
Il numeratore indica quante parti prendiamo, sono le parti $2\cdot 4=8$ in 
grigio. 

Il prodotto di due frazioni è una frazione che ha per numeratore il prodotto 
dei numeratori e per denominatore il prodotto dei denominatori.

\begin{center}
 % (c) 2012 Dimitrios Vrettos - d.vrettos@gmail.com
\begin{tikzpicture}
  \draw (0,0) rectangle  (10mm,10mm);
  \begin{scope}[font=\Huge]
    \node at  (5mm,5mm) {$\cdot$};
  \end{scope}
  \begin{scope}[font=\small]
    \node (a) at (-15mm,10mm) {$\displaystyle\frac{m}{n}$};
    \node (b) at (-15mm,0mm) {$\displaystyle\frac{p}{q}$};
    \node (c) at (25mm,5mm) {$\displaystyle\frac{mp}{nq}$};
  \end{scope}
  \node (entrata) at (0,5mm) {};
  \node (uscita) at (10mm,5mm) {};
  \begin{scope}[->,concept connection,blue,thin]
    \draw (a) -- (entrata);
    \draw (b) -- (entrata);
    \draw (uscita) -- (c);
  \end{scope}
\end{tikzpicture}

\end{center}

% \ovalbox{\risolvii \ref{ese:3.48}, \ref{ese:3.49}, \ref{ese:3.50}, \ref{ese:3.51}}

\subsection{Operazione inversa e aritmetica dell'orologio}

La divisione è l'operazione inversa della moltiplicazione. 
Ma cosa significa operazione inversa?
Una operazione può essere interpretata come qualsiasi azione che provoca 
un cambiamento di stato.

Consideriamo come esempio l'addizione nell'orologio che segna le ore 
dodici~$(12 =~0)$.
Addizionare significa spostare le lancette in avanti di un determinato 
numero di ore.
Si riporta la tabella dell'addizione dell'orologio.

Consideriamo l'addizione~$9+7=4$.
Il primo elemento~9 può essere interpretato come stato iniziale,~$+~7$ 
come operatore formato dall'operazione 
<<spostare le lancette avanti di\ldots>> e dall'argomento~7; 
il risultato~4 è lo stato finale.

Si indica come operazione inversa quella operazione che applicata allo stato 
finale con argomento uguale a quello precedente dell'operazione diretta, 
riporta allo stato iniziale.

Notiamo che anche nella matematica dell'orologio l'addizione gode della 
proprietà commutativa e associativa, ha l'elemento neutro che è~0, 
ogni numero ha l'inverso.

\newpage

\begin{center}
 % (c) 2012 Dimitrios Vrettos - d.vrettos@gmail.com
\begin{tikzpicture}[font=\small,x=6mm,y=6mm]
  % Aritmetica
  \node (o) at (0,0) {$+$}; % segno
  \foreach \x in {0,1,...,11} { 
    \node (a) at (\x+1,0) {\x}; % riga esterna
    \node (b) at (0,-\x-1) {\x}; % col. esterna
  }
  \draw (-3mm,-3mm)--(74mm,-3mm); % linea orizzontale
  \draw (3mm,3mm)--(3mm,-74mm); % linea verticale
  
  \foreach \x[count=\xi] in {0,...,11} 
    \foreach \y [count=\yi]in {\x,...,11}% 
      \node () at (\xi,-\yi) {\y};% Compila fino a 11, partendo da 0 e aumentando di 1 ad ogni nuova riga

  \foreach \x[count=\xii] in {0,...,10} 
    \foreach \y [count=\yii]in {\x,...,0}
      \node () at (\xii+1,\yii-13) {\y};% Compila fino a 10, partendo da 0 e aumentando di 1 fino alla fine di ogni nuova riga

  %Orologio
  \begin{scope}[xshift=10.9cm, yshift=-2.2cm] % spostamento
    \draw [double] (0,0) circle (2cm); % stile e dimensione
    \foreach \angle/ \label in {0/3, 30/2, 60/1, 90/12, 120/11, 150/10, 180/9,210/8, 240/7, 270/6, 300/5, 330/4}{% posizione ore
      \draw[line width=1pt] (\angle:1.8cm) -- (\angle:1.99cm); % lineette delle ore
      \draw (\angle:1.5cm) node{\textsf{\label}};
    }
    \begin{scope}[->,line width=1pt] % lancette
      \draw (0,0) -- (330:1cm); 
      \draw [dashed](0,0) -- (180:1cm);
    \end{scope}
    \path [decorate,decoration={text along path,text={indietro di 7}}](3.8,-2.2) arc  (-30:180:4.4); % testo curvato
    \path [decorate,decoration={text along path,text={avanti di 7}}](-3.9,0) arc  (180:-30:3.9); % testo curvato
    \draw[blue,<-<](-4.5,0) arc (180:-30:4.5);
    \draw[red,>->,dashed](-3.7,0) arc (180:-30:3.7);
  \end{scope}
  % Diagramma
  \begin{scope}[xshift=10.9cm,yshift=-4.7cm,->]
    \node at (0,0) {Operatore};
    \node at (-4,0) {Inizio};
    \node at (4,0) {Fine};

    \node [circle,fill=gray!40](nove) at (-4,-3) {9};
    \node[ellipse,fill=gray!20](avanti) at (0,-1.5)  {avanti di 7};
    \node[ellipse,fill=gray!20](indietro) at (0,-4.5)  {indietro di 7};
    \node [circle,fill=gray!40](quattro) at (4,-3) {4};

    \draw[red] (nove) ..controls +(up:10mm) and  +(left:15mm).. (avanti);
    \draw[red] (avanti) ..controls +(right:15mm) and  +(up:10mm).. (quattro);
    \draw[blue] (quattro) ..controls +(down:10mm) and  +(right:20mm).. (indietro);
    \draw[blue] (indietro) ..controls +(left:20mm) and  +(down:10mm).. (nove);
  \end{scope}
\end{tikzpicture}

\end{center}

\begin{multicols}{2}
\begin{itemize*}
\item L'inverso di~0 è~0 perché~$0+0=0$
\item l'inverso di~1 è~11 perché~$1+11=0$
\item l'inverso di~2 è~10 perché~$2+10=0$
\item l'inverso di~3 è~9 perché~$3+9=0$
\item l'inverso di~4 è~8 perché~$4+8=0$
\item l'inverso di~5 è~7 perché~$5+7=0$.
\end{itemize*}
\end{multicols}

L'elemento inverso è molto importante in quanto ci permette di sostituire 
l'operazione inversa, con l'operazione diretta che ha come argomento 
l'elemento inverso dell'argomento dell'operazione diretta.
%\newpage
\begin{center}
 % (c) 2012 Dimitrios Vrettos - d.vrettos@gmail.com
\begin{tikzpicture}[font=\small,x=6mm,y=6mm]
  %Orologio
  \begin{scope}[xshift=7cm,yshift=-18mm]
    \draw [double] (0,0) circle (2cm);
    \foreach \angle/ \label in{0/3, 30/2, 60/1, 90/12, 120/11, 150/10, 180/9,210/8, 240/7, 270/6, 300/5, 330/4}
    {
      \draw[line width=1pt] (\angle:1.8cm) -- (\angle:1.99cm);
      \draw (\angle:1.5cm) node{\textsf{\label}};
    }
    \begin{scope}[->,line width=1pt]
      \draw (0,0) -- (330:1cm); %
      \draw [dashed](0,0) -- (180:1cm);
    \end{scope}

    \path [decorate,decoration={text along path,text={avanti di 5}}](3.3,-2) arc  (330:180:3.7);
    \path [decorate,decoration={text along path,text={avanti di 7}}](-3.9,0) arc  (180:-30:3.9);
    \draw[blue,<-<](-3.7,0) arc (180:330:3.7);
    \draw[red,>->,dashed](-3.7,0) arc (180:-30:3.7);
  \end{scope}%
  % Diagramma
  \begin{scope}[->]
    \node at (0,0) {Operatore};
    \node at (-4,0) {Inizio};
    \node at (4,0) {Fine};

    \node [circle,fill=gray!40](nove) at (-4,-3) {9};
    \node[ellipse,fill=gray!20](avanti) at (0,-1.5)  {avanti di 7};
    \node[ellipse,fill=gray!20](indietro) at (0,-4.5)  {avanti di 5};
    \node [circle,fill=gray!40](quattro) at (4,-3) {4};

    \draw[red] (nove) ..controls +(up:10mm) and  +(left:15mm).. (avanti);
    \draw[red] (avanti) ..controls +(right:15mm) and  +(up:10mm).. (quattro);
    \draw[blue] (quattro) ..controls +(down:10mm) and  +(right:20mm).. (indietro);
    \draw[blue] (indietro) ..controls +(left:20mm) and  +(down:10mm).. (nove);
  \end{scope}
\end{tikzpicture}

\end{center}

Così per tornare allo stato iniziale invece di operare con portare 
indietro le lancette di~7, otteniamo lo stesso risultato portando avanti 
le lancette di~5 che è appunto l'inverso di~7.

\subsection{Divisione}

La divisione è l'operazione inversa della moltiplicazione.
Dato che nell'insieme dei numeri razionali esiste sempre l'inverso di una 
frazione rispetto alla moltiplicazione, esclusa la frazione zero, si può 
sempre eseguire la divisione di due qualsiasi frazioni.

\begin{center}
 % (c) 2012 Dimitrios Vrettos - d.vrettos@gmail.com
\begin{tikzpicture}
  \draw (0,0) rectangle  (10mm,10mm);
  \draw (40mm,0) rectangle  (50mm,10mm);

  \begin{scope}[font=\Huge]
    \node at  (5mm,5mm) {$:$};	
	\node at(15mm,5mm) {$=$};
	\node at(45mm,5mm) {$\cdot$};
  \end{scope}
  \begin{scope}[font=\small]
    \node (a) at (-15mm,10mm) {$\displaystyle\frac{m}{n}$};
    \node (b) at (-15mm,0mm) {$\displaystyle\frac{p}{q}$};
    \node (c) at (65mm,5mm) {$\displaystyle\frac{mq}{np}$};
    \node (ai) at (25mm,10mm) {$\displaystyle\frac{m}{n}$};
    \node (bi) at (25mm,0mm) {$\displaystyle\frac{q}{p}$};
  \end{scope}
  \node (entrata) at (0,5mm) {};
  \node (entratai) at (40mm,5mm) {};
  \node (uscita) at (50mm,5mm) {};
  \begin{scope}[->,concept connection,blue,thin]
    \draw (a) -- (entrata);
    \draw (b) -- (entrata);
    \draw (ai) -- (entratai);
    \draw (bi) -- (entratai);
     \draw (uscita) -- (c);
  \end{scope}
\end{tikzpicture}

\end{center}

\[\frac{m}{n}:\frac{p}{q}=\frac{m}{n}\cdot\frac{q}{p}=\frac{mq}{np}.\]

Il quoziente di due frazioni è la frazione che si ottiene moltiplicando la 
prima frazione per l'inverso della seconda frazione.

\begin{exrig}
 \begin{esempio}
Quoziente di due frazioni.
\begin{itemize*}
\item $\displaystyle{\frac{2}{3}:\frac{7}{4}}$.
\end{itemize*}
Il reciproco di~$\frac{7}{4}$ è~$\frac{4}{7}$. Pertanto
\[\frac{2}{3}:\frac{7}{4}\rightarrow\frac{2}{3}\cdot\frac{4}{7}=\frac{8}{21}.\]

\begin{itemize*}
	\item $\displaystyle{-\frac{2}{3}:\bigg(-\frac{3}{4}\bigg)}$.
\end{itemize*}
Il reciproco di~$-\frac{3}{4}$ è~$-\frac{4}{3}$. Pertanto
\[-\frac{2}{3}:\bigg(-\frac{3}{4}\bigg)\rightarrow-\frac{2}{3}\cdot
\bigg(-\frac{4}{3}\bigg)=+\frac{8}{9}.\]

\begin{itemize*}
 \item $\displaystyle{\frac{2}{3}:0}$.
\end{itemize*}
 Il reciproco di~0 non esiste, quindi la divisione non è eseguibile.

\begin{itemize*}
 \item $0:\dfrac{2}{3}$.
\end{itemize*}
Il reciproco di~$\frac{2}{3}$ è~$\frac{3}{2}$. Pertanto
\[0:\frac{2}{3}\rightarrow~0\cdot \frac{3}{2}=0.\]
\end{esempio}
\end{exrig}

% \ovalbox{\risolvii \ref{ese:3.52}, \ref{ese:3.53}, \ref{ese:3.54}, \ref{ese:3.55}}

\section{Potenza di una frazione}
\label{sec:03_potenza}

Come per ogni numero, anche per le frazioni, la potenza di una frazione non è
altro che un prodotto di tante frazioni identiche alla frazione data quanto è 
il valore dell'esponente, pertanto si trova elevando il numeratore e il 
denominatore della frazione all'esponente della potenza.
\[\bigg(\frac{a}{b}\bigg)^n=\underbrace{\frac{a}{b}\cdot\frac{a}{b}\cdot
\frac{a}{b}\cdot\ldots\cdot\frac{a}{b}}_%
{n\text{ volte}}=\frac{a^n}{b^n}.\]
\newpage
\begin{exrig}
\begin{esempio}
Potenza di frazioni.
\begin{multicols}{3}
\begin{itemize*}
\item $\displaystyle{\bigg(-\frac{2}{3}\bigg)^3=-\frac{8}{27}}$
\item $\displaystyle{-\frac{2^3}{3}=-\frac{8}{3}}$
\item $\displaystyle{\bigg(-\frac{2}{3}\bigg)^2=+\frac{4}{9}}$.
\end{itemize*}
\end{multicols}
\end{esempio}
\end{exrig}

\subsection{Potenza con esponente uguale a zero}
La definizione di potenza si estende anche al caso in cui l'esponente è zero.

Consideriamo l'esempio della divisione di due potenze con la stessa base e con 
lo stesso esponente:
\begin{itemize*}
\item $a^n:a^n=1$, la divisione di due numeri uguali è~1;
\item $a^n:a^n=a^0$, applicando le proprietà delle potenze.
\end{itemize*}

Possiamo allora concludere che per ogni frazione o numero razionale~$a$ 
diverso da zero~$a^0=1$. Non è invece possibile la potenza~$0^0$.

\subsection{Potenza con esponente un numero intero negativo}

La definizione di potenza si può estendere anche al caso in cui l'esponente 
sia uguale a un numero intero negativo:
\[a^{-n}=a^0:a^n=1:a^n=\frac{1}{a^n}=\frac{1^n}{a^n}=\bigg(\frac{1}{a}\bigg)^n.\]

Si può definire allora per ogni numero razionale diverso da zero
\[a^{-n}=\bigg(\frac{1}{a}\bigg)^n.\]

La potenza di un numero diverso da zero elevato a un esponente intero negativo 
è uguale a una potenza che ha per base il reciproco della base rispetto
alla moltiplicazione e per esponente l'opposto dell'esponente rispetto 
all'addizione.

Non è definita invece la potenza con esponente negativo di~0. Il numero~0 
infatti non ha il reciproco. 
Pertanto,~$0^{-n}$ è una scrittura priva di significato.

% \vspazio\ovalbox{\risolvii \ref{ese:3.56}, \ref{ese:3.57}, \ref{ese:3.58}, \ref{ese:3.59}, \ref{ese:3.60}}

\section{Notazione scientifica e ordine di grandezza}
\label{sec:03_ordinedigrandezza}

Le discipline scientifiche quali la fisica, la biologia, l'astronomia etc,
si trovano spesso a doversi confrontare con misurazioni di grandezze espresse 
da numeri molto grandi. Per esempio:

\begin{itemize*}
\item il raggio della Terra è circa~$6\,400\,000\unit{m}$
\item la velocità della luce nel vuoto è~$299\,790\,000\unit{m/s}$
\item un globulo rosso ha il diametro di~$0,000007\unit{m}$.
\end{itemize*}

I primi due numeri sono `molto grandi', mentre l'ultimo è `molto piccolo' e 
operare con numeri simili, non è affatto semplice.

Per renderci conto di ciò, consideriamo un rettangolo di dimensioni
$b =~0,00000006\unit{m}$ e~$h =~0,0000002\unit{m}$ e calcoliamone l'area:
\[A=b\cdot h=0,00000006\cdot0,0000002=0,000000000000012.\]

\begin{center}
% (c) 2012 Dimitrios Vrettos - d.vrettos@gmail.com
\begin{tikzpicture}[x=10mm,y=10mm,font=\small]

  \fill[fill=gray!10] (1,1) rectangle (7,3);
  \draw[style=help lines,step=5mm,dashed,black!20] (0.2,0.2) grid (7.8,3.8);
  \draw (1,1) rectangle (7,3);
  \node (b) at (4,.7) {$b$};
  \node (h) at (7.2,2) {$h$};
  \node [label={[name=label node]below left:$A$}] at (1,1) {};
  \node [label={[name=label node]above left:$B$}] at (1,3) {};
  \node [label={[name=label node]above right:$C$}] at (7,3) {};
  \node [label={[name=label node]below right:$D$}] at (7,1) {};

  \foreach \x in {1,7}
    \foreach \y in {1,3}
      \filldraw[fill=black, draw=black]  (\x,\y) circle (1pt);

\end{tikzpicture}

\end{center}

Come si può notare, per scrivere il risultato di un'operazione tra due numeri 
in questo caso `molto piccoli', è necessario fare particolare attenzione in 
quanto, per l'eccessiva quantità di cifre decimali, è facile commettere degli 
errori.

Per risolvere questo problema, si preferisce utilizzare una scrittura compatta 
che permette di scrivere questo tipo di numeri in forma più agevole.
Una tale scrittura prende il nome di \emph{notazione scientifica}.

\begin{definizione}
Un numero~$\alpha$ è scritto in \emph{notazione scientifica} se
si presenta nella forma:
\[\alpha=k\cdot10^n,\] 
dove~$k$ è un numero decimale maggiore o uguale a~$1$ e minore di~$10$ e~$n$ 
è un numero intero.
\end{definizione}

\begin{exrig}
 \begin{esempio}
I numeri~$3,5\cdot10^7$ e~$8,9\cdot10^{-5}$ sono scritti in notazione 
scientifica, mentre i numeri~$0,5\cdot10^3$ e~$10,3\cdot 10^{-8}$ non sono 
scritti in notazione scientifica in quanto il numero davanti alla potenza 
di~10 nel primo caso è~0,5 che è minore di~1, nel secondo caso è~10,3 che è 
maggiore di~10.
 \end{esempio}
\end{exrig}

\subsection{Come trasformare un numero in notazione scientifica}
Consideriamo la misura del diametro del globulo rosso, 
ovvero~$0,000007\unit{m}$.
Per esprimere questa misura in notazione scientifica basta considerare la sua 
frazione generatrice, ovvero:
\[0,000007\unit{m}=7\cdot\frac{1}{1000000}\unit{m}=7\cdot10^{-6}\unit{m}.\]

Allo stesso modo il numero~0,000000026 viene scritto in notazione scientifica 
come segue:
\[0,000000026=2,6\cdot\frac{1}{100000000}=2,6\cdot\frac{1}{10^8}=2,6\cdot10^{-8}.\]
Si osservi che in questo secondo caso abbiamo preso in considerazione il 
valore~2,6 anziché~26, in quanto il numero~$k$ deve essere minore di~10.

Consideriamo ora la misura del raggio della Terra, ovvero~$6\,400\,000\unit{m}$,
la sua espressione in notazione scientifica sarà:~$6,4\cdot10^6$.

Allo stesso modo il numero~$340\,000\,000\,000$ viene scritto in notazione 
scientifica~$3,4\cdot10^{11}$.
Si osservi che in questo secondo caso abbiamo preso in considerazione il 
valore~3,4 anziché~34, in quanto, come si è già detto, il numero~$k$ deve 
essere minore di~10.

\osservazione A numeri `piccoli', corrisponde una potenza di dieci con 
esponente negativo; 
a numeri `grandi', corrisponde una potenza di dieci con esponente positivo.

\begin{procedura}
Scrivere un numero decimale in notazione scientifica:
\begin{enumeratea}
 \item spostare la virgola di tanti posti in modo da avere una sola cifra
  diversa da zero a sinistra;
 \item scrivere la moltiplicazione tra il numero ottenuto al passo precedente
  e dieci elevato ad un esponente pari al numero si spostamenti della virgola 
  effettuati se la virgola è stata spostata verso sinistra o elevato al suo 
  opposto se la virgola è stata spostata verso destra.
\end{enumeratea}
\end{procedura}

\begin{exrig}
 \begin{esempio}
 Scrivi 348\,000\,000\,000\,000 in notazione scientifica.
 Per comodità riscrivo il numero evidenziando l'attuale 
 posizione della virgola:~348\,000\,000\,000\,000,0.
  \paragraph{Passo~$a$} Per ottenere un numero con una sola cifra diversa 
   da zero a sinistra della virgola devo spostare la virgola di~14 posti 
   verso sinistra;
  \paragraph{Passo~$b$} ora scrivo la moltiplicazione tra il numero 
   ottenuto:~$3,48$~e~10 elevato alla~14:~$3,48 \cdot 10^{14}$.
 \end{esempio}
 \begin{esempio}
  Scrivi 0,0000340 in notazione scientifica.
  \paragraph{Passo~$a$} Per ottenere un numero con una sola cifra diversa 
   da zero a sinistra della virgola devo spostare la virgola di~5 posti 
   verso destra;
  \paragraph{Passo~$b$} ora scrivo la moltiplicazione tra il numero 
   ottenuto:~$3,40$~e~10 elevato alla~$-5$:~$3,40 \cdot 10^{-5}$.
 \end{esempio}
 \begin{esempio}
  Riprendendo il problema della lamina rettangolare, 
  le sue dimensioni in notazione scientifica vengono scritte come:
  $b=6\cdot10^{-8}\unit{m}, h=2\cdot10^{-7}\unit{m}$.
  L'area sarà quindi:
  \begin{align*}
   A &=b\cdot h =\\
     &= 6\cdot10^{-8}\unit{m}\times 2\cdot10^{-7}\unit{m} =\\
     &=12\cdot 10^{-15}\unit{m}^2 =\\
     &=1,2\cdot 10^1\cdot 10^{-15}\unit{m}^2 =\\
     &=1,2\cdot 10^{-14}\unit{m}^2.
  \end{align*}
  Com'è possibile vedere, utilizzando le note proprietà delle potenze, 
  si riesce ad eseguire l'operazione in maniera molto agevole.
 \end{esempio}
 \begin{esempio}
  Trasforma in notazione scientifica e calcola
  $\displaystyle{\frac{3000:6\text{ milioni}}{5000\cdot0,000002}}$.
  \begin{align*}
   \frac{3000:6\text{ milioni}}{5000\cdot0,000002}&=\frac{3\cdot10^3:
     (6\cdot 10^6)}{5\cdot10^3\cdot(2\cdot10^{-6})}\\
   &=\frac{3:6\cdot10^{-3}}{5\cdot2\cdot 10^{-3}}\\
   &=\frac{0,5}{10}\cdot10^{-3+3}\\
   &= 0,05\cdot10^0\\
   &=0,05\\
   &=5\cdot10^{-2}.
  \end{align*}
 \end{esempio}
\end{exrig}

% \ovalbox{\risolvii \ref{ese:3.61}, \ref{ese:3.62}, \ref{ese:3.63}, 
% \ref{ese:3.64}, \ref{ese:3.65}, \ref{ese:3.66}, \ref{ese:3.67}, 
% \ref{ese:3.68}, \ref{ese:3.69}, \ref{ese:3.70}, \ref{ese:3.71}}

\subsection{Ordine di grandezza}

Spesso, nel trattare i numeri `molto grandi' o `molto piccoli', non è 
importante conoscere la misura con precisione, ma basta conoscere 
``quanto è grande'', cioè l'entità della sua grandezza. 
Per fare ciò si introduce il seguente concetto.

\begin{definizione}
Dato un numero scritto in forma scientifica, si definisce 
\emph{ordine di grandezza} (abbreviato con la sigla o.d.g.),
la potenza di~$10$.
\end{definizione}

\begin{procedura}
Determinare l'ordine di grandezza di un numero:
\begin{enumeratea}
\item scrivi il numero in notazione scientifica~${k\cdot10^n}$
\item l'ordine di grandezza è~$10^{n}$.
\end{enumeratea}
\end{procedura}

% \newpage

\begin{exrig}
 \begin{esempio}
Determinare l'ordine di grandezza dei numeri~0,000074 e~47000000000.

Scriviamo dapprima i numeri in notazione scientifica:
\[0,000074=7,4\cdot 10^{-5}\text{ e }47000000000=4,7\cdot 10^{10}.\]
L'o.d.g. del primo numero è~$10^{-5}$.
L'o.d.g del secondo numero è~$10^{10}$.
 \end{esempio}
\end{exrig}

% \ovalbox{\risolvii \ref{ese:3.72}, \ref{ese:3.73}, \ref{ese:3.74}}

\section{Problemi con le frazioni}
\label{sec:03_problemi}

\subsection{Problemi diretti}
Nei problemi diretti si conosce il valore di una grandezza e se ne deve 
calcolare la parte che corrisponde a una frazione. 
In questo caso basta moltiplicare la frazione per la grandezza intera.

\begin{exrig}
 \begin{esempio}
Una pasticceria produce~568 cornetti a settimana: i~$3/4$ sono alla 
crema,~$1/8$ sono al cioccolato e~$1/8$ alla marmellata. 
Quanti cornetti di ciascun tipo produce?

Per risolvere il problema occorre calcolare la parte che corrisponde a 
ciascuna frazione:

\begin{itemize*}
\item cornetti alla crema:~$\dfrac{3}{4}\cdot 568 =426$
\item cornetti al cioccolato:~$\dfrac{1}{8}\cdot 568 =71$
\item cornetti alla marmellata:~71.
\end{itemize*}
 \end{esempio}
\end{exrig}

\subsection{Problemi inversi}

Nei problemi inversi si conosce il valore numerico di una frazione di una 
certa grandezza si deve calcolare il valore dell'intera grandezza.
In questo caso occorre dividere il valore numerico dato per la frazione, 
si ottiene così l'intero.

\begin{exrig}
 \begin{esempio}
Mario ha speso~\officialeuro\ 21 che corrispondono ai~$3/5$ della somma 
che possedeva. Quanto possedeva?

In questo problema si sa che \officialeuro\ 21 corrispondono ai~$3/5$ della 
somma da cercare.
È sufficiente dividere~21 per la frazione:
\officialeuro\ 21~$\displaystyle:\frac{3}{5}=$ \officialeuro\ 21$\cdot\dfrac{5}{3}=$ \officialeuro\ 35.
 \end{esempio}

 \begin{esempio}
Giuseppe possiede \officialeuro\ 150. Se spende i~$3/5$ della somma posseduta e 
poi i~$2/3$ della somma rimanente, quanto gli rimane?

Per risolvere il problema si può procedere in più modi.

Calcoliamo prima i~$3/5$ di~150, cioè \officialeuro\ 150$\cdot\frac{3}{5}=$ 
\officialeuro\ 90.
Quindi la prima volta Giuseppe spende \officialeuro\ 90, perciò gliene 
rimangono~60. La seconda volta spende i~$2/3$ di \officialeuro\ 60, 
cioè \officialeuro\ 60$\frac{2}{3}=$ \officialeuro\ 40. In tutto ha speso
\officialeuro\ 90~$+$ \officialeuro\ 40~$=$ \officialeuro\ 130, gli 
rimangono \officialeuro\ 20.

Un altro modo per risolvere il problema è tenere conto che,
se la prima volta ha speso i~$\frac{3}{5}$ della somma che possedeva,
significa che gli rimane la frazione~$1-\frac{3}{5}=\frac{2}{5}$.
La seconda volta spende i~$\frac{2}{3}$ dei~$\frac{2}{5}$, cioè~$\frac{2}{3}
\cdot\frac{2}{5}=\frac{4}{15}$.
In tutto ha speso la frazione
\[\frac{3}{5}+\frac{4}{15}=\frac{3\cdot3+4}{15}=\frac{13}{15},\]
 gli rimane perciò la frazione~$\frac{2}{15}$, pertanto gli rimangono 
 \officialeuro\ 150$\cdot\frac{2}{15}=$ \officialeuro\ 20.
 \end{esempio}
\end{exrig}

% \ovalbox{\risolvii \ref{ese:3.75}, \ref{ese:3.76}, \ref{ese:3.77}, \ref{ese:3.78}}

\section{Le percentuali}
\label{sec:03_percentuali}

Avrai sentito parlare spesso che il prezzo di un oggetto è stato scontato 
del~10 per cento, oppure che un partito politico ha preso il~25 per cento di 
voti e altre espressioni simili che coinvolgono le percentuali.

Le percentuali sono un altro modo per scrivere le frazioni.

\begin{definizione}
Le \emph{percentuali} sono frazioni che hanno come denominatore~100 e come 
numeratore un numero intero o decimale.
\end{definizione}

La percentuale si indica con un numero intero o decimale seguita dal simbolo \%.
\[35\%=\frac{35}{100};\qquad7\%=\frac{7}{100};\qquad~12,5\%=\frac{12,5}{100}=
\frac{125}{1000}.\]

Per passare dalla scrittura percentuale alla scrittura decimale basta dividere 
per~100 il numero che esprime la percentuale:
\[35\%=\frac{35}{100}=0,35;\qquad7\%=\frac{7}{100}=0,07;\qquad12,5\%=
\frac{12,5}{100}=0,125.\]

Per passare dalla scrittura decimale alla scrittura in percentuale basta 
moltiplicare numeratore e denominatore per~100:
\[0,02=\frac{0,02}{1}=\frac{2}{100}=2\%;\qquad0,23=\frac{0,23}{1}=
\frac{23}{100}=23\%; \qquad1,21=\frac{1,21}{1}=\frac{121}{100}=121\%.\]

Per passare da una frazione alla percentuale conviene prima scrivere la 
frazione come numero decimale e poi da questo passare alla percentuale:
\[\frac{2}{3}=0,\overline{6}=\frac{0,\overline{6}}{1}=\frac{66,\overline{6}}{100}=66,
\overline{6}\%.\]

% \vspazio\ovalbox{\risolvii \ref{ese:3.79}, \ref{ese:3.80}, \ref{ese:3.81}, \ref{ese:3.82}}

\subsection{Problemi con le percentuali}

Per calcolare la percentuale di una grandezza è sufficiente moltiplicare il 
valore della grandezza per la percentuale espressa in frazione.

\begin{exrig}
 \begin{esempio}
In una scuola che ha~857 alunni ne sono stati promossi il~95\%. 
Quanti sono stati i promossi?

Per rispondere, si moltiplica il numero totale di alunni per la 
frazione~$95/100$.
Precisamente~$\frac{95}{100}\cdot857=814,15$.
Poiché il risultato non è un numero intero la percentuale è stata 
approssimata. Gli alunni promossi sono stati~814.
 \end{esempio}
\end{exrig}

A volte è nota una parte della grandezza e si vuole conoscere che percentuale è
la parte nota rispetto al totale. In questo caso occorre dividere la parte 
nota per l'intera grandezza, moltiplicare il risultato per~100 ed esprimere 
il numero in percentuale.

\begin{exrig}
\begin{esempio}
Di una scolaresca di~652 alunni ben~126 hanno avuto il debito in matematica.
Qual è la percentuale di alunni che hanno avuto il debito in matematica?

Per rispondere alla domanda eseguiamo i seguenti calcoli:
\[\frac{126}{652}\cdot100\%\approx0,19\cdot100\%=19\%.\]
\end{esempio}
\end{exrig}

\subsection{Problemi con gli sconti}

\begin{exrig}
 \begin{esempio}
Un pantalone costava \officialeuro\ 70 e viene venduto con il~$20\%$ di sconto, 
a quanto viene venduto?

Si tratta di calcolare prima lo sconto e po il prezzo scontato. 
Lo sconto è dato da

\begin{center}
$\displaystyle{20\%\cdot70}$ \officialeuro\ $=\dfrac{20}{100}\cdot70$ 
\officialeuro\ $=14$.
\end{center}

Il prezzo scontato è \officialeuro\ $70 -$ \officialeuro\ 14~$=$ 
\officialeuro\ 56.

In alternativa si può tenere conto che, se~$20\%$ esprime lo sconto, la parte 
rimanente, quella da pagare, è $100\%-20\%=80\%$. 
Quindi per calcolare quanto costano i pantaloni scontati si può calcolare

\begin{center}
$\displaystyle{80\%\cdot70}$ \officialeuro\ $=\dfrac{80}{100}\cdot70$ 
\officialeuro\ $=56$ \officialeuro.
 \end{center}
 \end{esempio}

 \begin{esempio}
Un paio di scarpe da \officialeuro\ 120 viene venduto scontato a 
\officialeuro\ 75 Qual è stata la percentuale di sconto praticato?

Per rispondere alla domanda, calcolo lo sconto \officialeuro\ 120~$-$ 
\officialeuro\ 75~$=$ \officialeuro\ 45.

Calcolo la percentuale che \officialeuro\ 45 rappresentano di 
\officialeuro\ 120, \[\frac{45}{120}\cdot100\%=0,375\cdot100\% =~37,5\%.\]
 \end{esempio}

\begin{esempio}
Mario ha trovato in un negozio il computer che stava cercando; per fortuna 
era scontato del~$15\%$, ha risparmiato cosi~120 euro. 
Quanto costa il computer di listino?

\officialeuro\ 120 corrispondono al~$15\%$ del prezzo di listino.
Per calcolare il prezzo di listino occorre dividere~120 per la frazione 
che corrisponde a~$15\%$.
\begin{center}
 $120:~15\% =~120:\dfrac{15}{100} =120\cdot\dfrac{100}{15}=$ \officialeuro\ 800.
\end{center}
\end{esempio}
\end{exrig}


% \ovalbox{\risolvii \ref{ese:3.83}, \ref{ese:3.84}, \ref{ese:3.85}, 
% \ref{ese:3.86}, \ref{ese:3.87},
% \ref{ese:3.88}, \ref{ese:3.89}, \ref{ese:3.90}, \ref{ese:3.91}, 
% \ref{ese:3.92}, \ref{ese:3.93},
% \ref{ese:3.94}, \ref{ese:3.95},}

% \vspazio\ovalbox{\ref{ese:3.96}, \ref{ese:3.97}, \ref{ese:3.98}, 
% \ref{ese:3.99}, \ref{ese:3.100}, \ref{ese:3.101},
% \ref{ese:3.102}, \ref{ese:3.103}, \ref{ese:3.104}, \ref{ese:3.105}, 
% \ref{ese:3.106}, \ref{ese:3.107}, \ref{ese:3.108},
% \ref{ese:3.109},}
% 
% \vspazio\ovalbox{\ref{ese:3.110}, \ref{ese:3.111}, \ref{ese:3.112}, 
% \ref{ese:3.113}, \ref{ese:3.114}}

\section{Proporzioni}
\label{sec:03_proporzioni}

\begin{definizione}
Il rapporto tra due numeri, di cui il secondo è diverso da zero, è il 
quoziente che si ottiene dividendo il primo numero per il secondo. 
Il primo numero si dice \emph{antecedente}, il secondo \emph{conseguente}.
\end{definizione}

\begin{center}
  % (c) 2012 Dimitrios Vrettos - d.vrettos@gmail.com
\begin{tikzpicture}[x=5mm,y=5mm]
  \node (a) at (-2,0) {$A$};
  \node (b) at (2,0) {$B$};
  \node (c) at (0,0) {$:$};
    \begin{scope}[font=\small]
    \node (ter) at (0,2) {termini};
    \node (ant) at (-4,-2) {antecedente};
    \node (con) at (4,-2) {conseguente};
    \end{scope}
    \begin{scope}[->]
      \begin{scope}[blue]
	\draw (ter) ..controls +(left:20mm) and +(left:10mm).. (a);
	\draw (ter) ..controls +(right:20mm) and +(right:10mm).. (b);
      \end{scope}
      \begin{scope}[red]
	\draw (ant) ..controls +(up:5mm) and +(down:10mm).. (a);
	\draw (con) ..controls +(up:5mm) and +(down:10mm).. (b);
      \end{scope}
    \end{scope}
\end{tikzpicture}

\end{center}

\begin{definizione}
 Una \emph{proporzione} è una uguaglianza tra due rapporti, del tipo
\[A: B = C: D,\]
che si legge \emph{$A$ sta a~$B$ come~$C$ sta a~$D$}, con~$B$ e~$D$ diversi da 
zero.
\end{definizione}

\begin{center}
  % (c) 2012 Dimitrios Vrettos - d.vrettos@gmail.com

\begin{tikzpicture}[x=5mm,y=5mm]
  \node (a) at (-3,0) {$A$};
  \node (b) at (-1,0) {$B$};
  \node (e) at (0,0) {$=$};
  \node (c) at (1,0) {$C$};
  \node (d) at (3,0) {$D$};
  \foreach \x in {-2,2}
    \node (f) at (\x,0) {$:$};
  \begin{scope}[font=\small]
    \node (est) at (0,3) {estremi};
    \node (med) at (0,2) {medi};
    \node (ant) at (-4,-2) {antecedenti};
    \node (con) at (4,-2) {conseguenti};
  \end{scope}
  \begin{scope}[->]
    \begin{scope}[blue]
      \draw (est) ..controls +(left:15mm) and +(up:6mm).. (a);
      \draw (est) ..controls +(right:15mm) and +(up:10mm).. (d);
    \end{scope}
    \begin{scope}[red]
      \draw (med) ..controls +(down:5mm) and +(up:8mm).. (b);
      \draw (med) ..controls +(down:5mm) and +(up:8mm).. (c);
    \end{scope}
    \begin{scope}[orange]
      \draw (ant)..controls +(up:2mm) and +(left:8mm).. (a);
      \draw (ant)..controls +(up:5mm) and +(down:10mm).. (c);
    \end{scope}
    \begin{scope}[green]
      \draw (con)..controls +(up:5mm) and +(down:8mm).. (b);
      \draw (con)..controls +(up:5mm) and +(right:8mm).. (d);
    \end{scope}
  \end{scope}
\end{tikzpicture}

\end{center}

\begin{exrig}
 \begin{esempio}
  $4:~2 =~12:~6$.

Formano una proporzione perché i due quozienti valgono entrambi~2.
 \end{esempio}

\begin{esempio}
$7:~14 =~16:~4$.

 \emph{Non} formano una proporzione perché il primo rapporto vale~0,5 mentre 
 il secondo rapporto vale~4.
\end{esempio}
\end{exrig}

Si dice anche che quattro numeri sono in proporzione se il rapporto tra i 
primi due è uguale al rapporto tra il terzo e il quarto.
 
\begin{proprieta}[Fondamentale delle proporzioni]
  In ogni proporzione il prodotto dei medi è uguale al prodotto degli estremi.
\[A:B=C:D\Rightarrow A\cdot D=B\cdot C.\]
\end{proprieta}

\begin{exrig}
 \begin{esempio}
$4:~6 =~6:~9$.

Il prodotto dei medi è~$6\cdot6=36$ e il prodotto degli estremi è~$4\cdot9=36$. 
Quindi è una proporzione.
 \end{esempio}

\begin{esempio}
$20:~30 =~30:~40$.

Il prodotto dei medi è~$30\cdot30=900$ il prodotto degli estremi 
è~$20\cdot40=800$. Quindi non è una proporzione.
\end{esempio}
\end{exrig}

% \begin{proprieta}[del permutare]
% Se in una proporzione scambiamo tra di loro i medi otteniamo ancora una 
% proporzione; in modo analogo otteniamo ancora una proporzione se scambiamo 
% tra di loro gli estremi, o ancora se scambiamo tra di loro sia i medi sia 
% gli estremi.
% \[A: B = C: D\Rightarrow A: C = B: D\Rightarrow D: B = C: A\Rightarrow 
% D: C = B: A.\]
% \end{proprieta}
% 
% \begin{exrig}
%  \begin{esempio}
% Data la proporzione~$12:~16 =~18:~24$ e scambiando tra di loro:
% \begin{itemize*}
%   \item i medi si ottiene la proporzione~$12:~18 =~16:~24$
%   \item gli estremi si ottiene la proporzione~$24:~16 =~18:~12$
%   \item sia i medi sia gli estremi si ottiene la 
%     proporzione~$24:~18 =~16:~12$.
% \end{itemize*}
% 
%  \end{esempio}
% \end{exrig}
% 
% 
% \begin{proprieta}[dell'invertire]
%   Se in una proporzione scambiamo ogni antecedente con il rispettivo 
%   conseguente
% otteniamo ancora una proporzione.
% \[A: B = C: D\Rightarrow B: A = D: C.\]
% \end{proprieta}
% 
% \begin{exrig}
%  \begin{esempio}
% Data la proporzione~$15:~9 =~5:~3$, applicando la proprietà dell'invertire
% otteniamo la proporzione~$9:~15 =~3:~5$.
%  \end{esempio}
% \end{exrig}
% 
% \begin{proprieta}[del comporre]
% In una proporzione la somma dei primi due termini sta al primo termine come 
% la somma del terzo e del quarto termine sta al terzo termine. Analogamente,
% la somma dei primi due termini sta al secondo termine come la somma del 
% terzo e del quarto termine sta al quarto termine.
% \[A: B = C: D\Rightarrow (A+B): A = (C+D): C.\]
% \[A: B = C: D\Rightarrow (A+B): B = (C+D): D.\]
% \end{proprieta}
% 
% \begin{exrig}
%  \begin{esempio}
% Data la proporzione~$16:~10 =~40:~25$, applicando la proprietà del comporre 
% si ottengono le proporzioni
% \[26:~16 =~65:~40,\qquad26:~10 =~65:~25.\]
%  \end{esempio}
% \end{exrig}
% 
% Analogamente alla proprietà del comporre si ha la seguente:
% 
% \begin{proprieta}[dello scomporre]
%   In una proporzione la differenza dei primi due termini sta al primo 
% termine come la differenza del
% terzo e del quarto termine sta al terzo termine. Analogamente, la differenza 
% dei primi due termini % sta al secondo termine come la differenza del terzo 
% e del quarto termine sta al quarto termine.
% \[A: B = C: D\Rightarrow (A-B): A = (C-D): C.\]
% \[A: B = C: D\Rightarrow (A-B): B = (C-D): D.\]
% \end{proprieta}
% 
% \begin{exrig}
%  \begin{esempio}
% Data la proporzione~$16:~10 =~40:~25$, applicando la proprietà dello scomporre 
% si ottengono le proporzioni
% \[6:~16 =~15:~40,\qquad6:~10 =~15:~25.\]
%  \end{esempio}
% \end{exrig}

\subsection{Calcolo di un medio o un estremo incognito}
Il medio incognito di una proporzione si calcola moltiplicando gli estremi e 
dividendo il risultato per l'altro medio:
\[a:b=x:d\Rightarrow x=\frac{a\cdot d}{b}.\]

L'estremo incognito di una proporzione si calcola moltiplicando i medi e 
dividendo il risultato per l'altro estremo:
\[x:b=c:d\Rightarrow x=\frac{b\cdot c}{d}.\]

\begin{exrig}\vspace{-1ex}
\begin{esempio}
Calcola il termine incognito di ciascuna proporzione.
\vspace{-1.3ex}\begin{itemize*}
  \item $5:7=20:x\Rightarrow x=\frac{7\cdot 20}{5}=28$
  \item $2:x=~3:16\Rightarrow x=\frac{2\cdot 16}{3}=\frac{32}{3}$
  \item $\frac{2}{3}:\frac{1}{2}=x:\frac{5}{6}%
\Rightarrow x=\frac{2}{3}\cdot\frac{5}{6}:\frac{1}{2}=\frac{2}{3}%
\cdot\frac{5}{6}\cdot\frac{2}{1}=\frac{10}{9}$.
\end{itemize*}
\end{esempio}%\vspace*{-1.7ex}
\end{exrig}

\begin{definizione}
  Una proporzione si dice \emph{continua} se ha i medi uguali.
\end{definizione}

Una proporzione continua è del tipo~$A: B = B: C$, per esempio
\[3:~9 =~9:~27,\qquad5:~10 =~10:~20,\qquad4:~16 =~16:~64.\]

\subsubsection*{Calcolo del medio in una proporzione continua}

In una proporzione continua il medio proporzionale incognito si ottiene 
moltiplicando gli estremi e calcolando la radice quadrata del prodotto ottenuto.
\[a:x=x:d\Rightarrow x=\sqrt{a\cdot d}.\]

\begin{exrig}
\begin{esempio}
Trovare il valore di x nella seguente proporzione continua~$36: x = x:~9$.

Svolgimento~$x=\sqrt{36\cdot 9}=18$.
\end{esempio}
\end{exrig}

% \subsubsection*{Calcolo di un termine incognito per mezzo delle proprietà del 
% comporre e dello scomporre}
% 
% \begin{exrig}
%   \begin{esempio}
%     $(11-x): x =~15:~5$.
% 
%   Applicando la proprietà del comporre si ha la proporzione
%   \begin{align*}
%   (11-x+x): x = (15+5):~5 &\Rightarrow~11:x=20:5\\
%   &\Rightarrow x=\frac{11\cdot5}{20}=\frac{11}{4}.
%   \end{align*}
%   \end{esempio}
% 
% \begin{esempio}
%  ~$\displaystyle{\bigg(\frac{1}{2}+x\bigg):\frac{5}{8}=x:5}$.
% 
%  Permutando i medi si ha~$\displaystyle{\bigg(\frac{1}{2}+x\bigg):x=
% \frac{5}{8}:5}$. Applicando la proprietà
% dello scomporre si ha:
% \begin{align*}
%   &\bigg(\frac{1}{2}+x-x\bigg):x=\bigg(\frac{5}{8}-5\bigg):5\\
% \Rightarrow &\frac{1}{2}:x=\frac{-35}{8}:5\\
% \Rightarrow &x=\frac{1}{2}\cdot5:\bigg(\frac{-35}{8}\bigg)=\frac{1}{2}\cdot5%
% \cdot\bigg(-\frac{8}{35}\bigg)=-\frac{4}{7}.
% \end{align*}
% \end{esempio}
% \end{exrig}

\subsection{Grandezze direttamente e inversamente proporzionali}

Si consideri il perimetro di un triangolo equilatero;
sappiamo che esso varia al variare della lunghezza del suo lato.
Se si indica con~$l$ la lunghezza del lato del triangolo, allora il perimetro 
è dato dalla relazione:~\[2p=3l.\]

È possibile notare che se raddoppia il lato, raddoppia anche il perimetro; 
se si triplica il lato, allora triplica anche il perimetro etc.

\begin{center}
 \begin{tabular*}{.9\textwidth}{l@{\extracolsep{\fill}}*{6}{c}}
\toprule
Lato~$l$&~0,5&~1&~1,5&~2,4&~3,1&~4,4\\
Perimetro & 1,5& 3& 4,5& 7,2& 9,3& 13,2\\
Rapporto~$\dfrac{2p}{l}$ &~3&~3&~3&~3&~3&~3\\
\bottomrule
\end{tabular*}
\end{center}

\begin{definizione}
  Due grandezze~$x$ e~$y$ si dicono \emph{direttamente proporzionali} se il 
  loro rapporto è costante, cioè \[\frac{y}{x}=k,\text{ con }k\neq~0.\]
\end{definizione}

In generale, da quest'ultima scrittura, possiamo dedurre che una 
proporzionalità diretta è espressa da una formula del tipo:
\[y=kx,\text{ con }k\neq 0.\]

Graficamente un tale tipo di proporzionalità è rappresentato da una retta che
passa per l'origine di un sistema di assi cartesiani ortogonali 
(figura~\ref{fig:3.2}).

\begin{inaccessibleblock}[Figura: TODO]
 \begin{figure}[h]
\begin{minipage}[t]{.48\textwidth}
  \centering% (c) 2012 Dimitrios Vrettos - d.vrettos@gmail.com
\begin{tikzpicture}[domain=0:4.4,x=10mm,y=5mm]
  \draw[->] (0,0) -- coordinate (x axis mid) (5.2,0); %asse x
  \draw[->] (0,0) -- coordinate (y axis mid) (0,14.5); % asse y
  % etichette
  \foreach \x in {0,1,...,5}
    \draw (\x,3pt) -- (\x,-3pt) node[anchor=north] {\x};
  \foreach \y in {1,2,...,14}
    \draw (3pt,\y) -- (-3pt,\y) node[anchor=east] {\y}; 
  \node[below=0.5cm] at (x axis mid) {lato $l$};
  \node[rotate=90,above left=1cm] at (y axis mid) {perimetro $p$};
  %grafico
  \draw[color=blue] plot  (\x,{3*\x});
  \begin{scope}[dotted]
    \foreach \x in {1,2,3,4}{
      \draw (\x,0) -- (\x,3*\x);
      \draw (0,3*\x) -- (\x,3*\x);
    }
  \end{scope}
\end{tikzpicture}

  \caption{Proporzionalità diretta.}\label{fig:3.2}  
\end{minipage}\hfil
 \begin{minipage}[t]{.48\textwidth}  
 \centering% (c) 2012 Dimitrios Vrettos - d.vrettos@gmail.com
\begin{tikzpicture}[x=10mm,y=10mm, smooth]
  \draw[->] (0,0) -- coordinate (x axis mid) (5.2,0);
  \draw[->] (0,0) -- coordinate (y axis mid) (0,4.5);
  \foreach \x in {0,1,...,5}
    \draw (\x,3pt) -- (\x,-3pt) node[anchor=north] {\x};
  \foreach \y in {1,2,...,4.5}
    \draw (3pt,\y) -- (-3pt,\y) node[anchor=east] {\y}; 
  %labels      
  \node[below=0.5cm] at (x axis mid) {volume $V$};
  \node[rotate=90,above left=1cm] at (y axis mid) {pressione $P$};
  \draw[domain=.5:4.4,color=blue] plot  (\x,{2/\x});
  \begin{scope}[dotted]
    \foreach \x in {.5,1,2,3,4.4}{
      \draw (\x,0) -- (\x,2/\x);
      \draw (0,2/\x) -- (\x,2/\x);
    }
  \end{scope}
\end{tikzpicture}

 \caption{Proporzionalità inversa.}\label{fig:3.3}
 \end{minipage}
\end{figure}
\end{inaccessibleblock}

Esaminiamo ora un altro esempio. Se quando vai a fare benzina allo scooter 
chiedi ogni volta \officialeuro\ 10 di benzina, noterai che se aumenta il 
prezzo della benzina diminuirà la quantità di carburante che ricevi e 
viceversa se diminuisce il prezzo aumenterà la quantità di carburante che 
ricevi. Ciò che rimane costante è il prodotto tra il prezzo della benzina e 
la quantità di benzina ricevuta che deve essere sempre \officialeuro\ 10.

\begin{center}
 \begin{tabular*}{.9\textwidth}{l@{\extracolsep{\fill}}*{4}{c}}
\toprule
Prezzo benzina al litro~$p$ (\officialeuro\ ) &1,126 &1,156 &1,212 &1,248\\
Benzina ricevuta~$b$ (\unit{l}) &8,881 &8,650 &8,251 &8,013\\
Costo~$c=p\cdot b $ (\officialeuro\ ) &10,00 &10,00 &10,00 &10,00\\
\bottomrule
\end{tabular*}
\end{center}

\begin{definizione}
  Due grandezze~$x$ e~$y$ si dicono \emph{inversamente proporzionali} se il 
  loro prodotto è costante,
cioè se:
\[x\cdot y=k,\text{ con }k\neq 0.\]
\end{definizione}

In generale, da quest'ultima scrittura, possiamo dedurre che una 
proporzionalità diretta è espressa da una formula del tipo:
\[y=\frac{k}{x},\text{ con }k\neq~0.\]

Graficamente un tale tipo di proporzionalità è rappresentato da un ramo 
d'iperbole equilatera in un sistema di assi cartesiani ortogonali 
(figura~\ref{fig:3.3}).

% \vspazio\ovalbox{\risolvii \ref{ese:3.115}, \ref{ese:3.116}, 
% \ref{ese:3.117}, \ref{ese:3.118}, \ref{ese:3.119}, \ref{ese:3.120},
% \ref{ese:3.121}, \ref{ese:3.122}, \ref{ese:3.123}, \ref{ese:3.124}, 
% \ref{ese:3.125}}
% 
% \vspazio\ovalbox{\ref{ese:3.126}, \ref{ese:3.127}}

\clearpage

\section{Espressioni con le frazioni}
\label{sec:03_espressioni}

% \begin{esempio}
% %\begin{problema}
%   Calcola il valore della seguente espressione.
% 
% \begin{figure}[h]
%  \begin{minipage}[]{.60\textwidth}
% 
% $\bigg\lbrace
%   \frac{3}{20}\cdot\bigg[
%     \bigg(\frac{4}{9}-\frac{1}{3}\bigg):5 + 
%     \bigg(\frac{3}{7}-\frac{2}{5}\bigg):
%     \frac{1}{14}+\frac{1}{5}\cdot\frac{1}{9}\bigg]+
%   \frac{2}{15}\bigg\rbrace :
% 2=$
% \vspace{2cm}
% 
% $=\bigg\lbrace
%   \frac{3}{20}\cdot\bigg[
%     \frac{4 - 3}{9} \cdot \frac{1}{5} + 
%     \frac{15-14}{35}\cdot\frac{14}{1} +
%     \frac{1}{45}\bigg] +
%   \frac{2}{15}\bigg\rbrace \cdot 
% \frac{1}{2}=$
% 
% $=\bigg\lbrace
%   \frac{3}{20}\cdot\bigg[
%     \frac{1}{9} \cdot \frac{1}{5} + 
%     \frac{1}{35}\cdot\frac{14}{1} +
%     \frac{1}{45}\bigg] +
%   \frac{2}{15}\bigg\rbrace \cdot 
% \frac{1}{2}=$
% 
% $=\bigg\lbrace
%   \frac{3}{20}\cdot\bigg[
%     \frac{1}{45} + 
%     \frac{2}{5} +
%     \frac{1}{45}\bigg] +
%   \frac{2}{15}\bigg\rbrace \cdot 
% \frac{1}{2}=$
% 
% $=\bigg\lbrace
%   \frac{3}{20}\cdot\bigg[
%     \frac{1+18+1}{45}\bigg] +
% \frac{2}{15}\bigg\rbrace \cdot \frac{1}{2}=$
% 
% $=\bigg\lbrace
%   \frac{3}{20}\cdot \frac{20}{45} + \frac{2}{15}\bigg\rbrace \cdot 
% \frac{1}{2}=$
% 
% $=\bigg\lbrace
%   \frac{1}{15} + \frac{2}{15}\bigg\rbrace \cdot 
% \frac{1}{2}=$
% 
% $=\frac{3}{15} \cdot \frac{1}{2} = \frac{1}{6}$
% 
%  \end{minipage}
%  \begin{minipage}[]{.35\textwidth}
% \begin{enumerate}
%  \item 
% Eseguo le operazioni contenute nelle parentesi più interne. 
% La linea di frazione equivale ad una coppia 
% di parentesi per cui le parentesi tonde non sono più necessarie. 
% Trasformo le divisioni in moltiplicazioni per il reciproco del divisore.
%  \item
% Eseguo le addizioni presenti nei numeratori di due frazioni.
%  \item
% Eseguo le moltiplicazioni tra frazioni presenti nella parentesi quadra.
%  \item
% Addiziono le frazioni presenti nella quadra.
%  \item
% Eseguo la moltiplicazione nella parentesi graffa,
%  \item
% Eseguo l'addizione nella graffa.
%  \item
% Eseguo l'ultima moltiplicazione.
% \end{enumerate}
%  \end{minipage}
% \end{figure}
% 
% \end{esempio}
\begin{esempio}
%\begin{problema}
  Calcola il valore della seguente espressione.

\[\bigg\lbrace
  \frac{3}{20}\cdot\bigg[
    \bigg(\frac{4}{9}-\frac{1}{3}\bigg):5 + 
    \bigg(\frac{3}{7}-\frac{2}{5}\bigg):
    \frac{1}{14}+\frac{1}{5}\cdot\frac{1}{9}\bigg]+
  \frac{2}{15}\bigg\rbrace :
2=\]
\begin{enumerate}
 \item 
Eseguo le operazioni contenute nelle parentesi più interne. 
La linea di frazione equivale ad una coppia 
di parentesi per cui le parentesi tonde non sono più necessarie. 
Trasformo le divisioni in moltiplicazioni per il reciproco del divisore.
\[=\bigg\lbrace
  \frac{3}{20}\cdot\bigg[
    \frac{4 - 3}{9} \cdot \frac{1}{5} + 
    \frac{15-14}{35}\cdot\frac{14}{1} +
    \frac{1}{45}\bigg] +
  \frac{2}{15}\bigg\rbrace \cdot 
\frac{1}{2}=\]
 \item
Eseguo le addizioni presenti nei numeratori di due frazioni.
\[=\bigg\lbrace
  \frac{3}{20}\cdot\bigg[
    \frac{1}{9} \cdot \frac{1}{5} + 
    \frac{1}{35}\cdot\frac{14}{1} +
    \frac{1}{45}\bigg] +
  \frac{2}{15}\bigg\rbrace \cdot 
\frac{1}{2}=\]
 \item
Eseguo le moltiplicazioni tra frazioni presenti nella parentesi quadra.
\[=\bigg\lbrace
  \frac{3}{20}\cdot\bigg[
    \frac{1}{45} + 
    \frac{2}{5} +
    \frac{1}{45}\bigg] +
  \frac{2}{15}\bigg\rbrace \cdot 
\frac{1}{2}=\]
 \item
Addiziono le frazioni presenti nella quadra.
\[=\bigg\lbrace
  \frac{3}{20}\cdot\bigg[
    \frac{1+18+1}{45}\bigg] +
\frac{2}{15}\bigg\rbrace \cdot \frac{1}{2}=\]
 \item
Eseguo la moltiplicazione nella parentesi graffa,
\[=\bigg\lbrace
  \frac{3}{20}\cdot \frac{20}{45} + \frac{2}{15}\bigg\rbrace \cdot 
\frac{1}{2}=\]
\[=\bigg\lbrace
  \frac{1}{15} + \frac{2}{15}\bigg\rbrace \cdot 
\frac{1}{2}=\]
 \item
Eseguo l'addizione nella graffa.
\[=\frac{3}{15} \cdot \frac{1}{2} = \frac{1}{6}\]
 \item
Eseguo l'ultima moltiplicazione.
\[= \frac{1}{6}\]

\end{enumerate}

\end{esempio}
% \newpage

Quando si deve addizionarre un numero intero ad una frazione si può evitare 
di usare un passaggio per scrivere esplicitamente il denominatore comune: 
dato che c'è un unico denominatore è ovvio che quello \emph{è} il 
denominatore comune. È possibile ottenere la somma di un intero e una 
frazione con un unico passaggio moltiplicando il numero intero per il denominatore 
e aggiungendolo al numeroatore della frazione.

\clearpage

\begin{esempio}
Esegui i seguenti calcoli:

\begin{itemize*}
 \item 
$3+\frac{5}{4} = \dots$\qquad
<<tre per quattro -> dodici, più cinque -> diciassette, quarti>>
 \item
$5-\frac{1}{2} = \dots$\qquad
<<cinque per due -> dieci, meno uno -> nove, mezzi>>
 \item
$\frac{8}{3} - 2 = \dots$\qquad
<<due per tre -> sei, otto meno sei -> due, terzi>>
\end{itemize*}

\end{esempio}

Applicando questo metodo si può risparmiare qualche passaggio.

\begin{esempio}
%\begin{problema}
  Calcola il valore della seguente espressione.
  
% \[\bigg[\frac{13}{5}:\bigg(3+\frac{9}{10}\bigg)+\frac{7}{8}+
% \bigg(\frac{13}{4}-2\bigg)\cdot\frac{4}{15}-\frac{7}{8}%
% \bigg]\cdot\frac{11}{3}:\bigg(6-\frac{1}{2}\bigg).\]
%\end{problema}
%\begin{soluzione}
\begin{flalign*}
&\bigg[
  \frac{13}{5}:\bigg(3+\frac{9}{10}\bigg)+\frac{7}{8}+
  \bigg(\frac{13}{4}-2\bigg)\cdot\frac{4}{15}-\frac{7}{8}%
\bigg]\cdot
\frac{11}{3}:\bigg(6-\frac{1}{2}\bigg)=&\\
&=\bigg[
  \frac{13}{5}:\frac{39}{10}+\frac{7}{8}+
  \frac{5}{4}\cdot\frac{4}{15}-\frac{7}{8}
\bigg]\cdot
\frac{11}{3}:\frac{11}{2}=\\
&=\bigg[
  \frac{13}{5}\cdot\frac{10}{39}+\frac{7}{8}+
  \frac{5}{4}\cdot\frac{4}{15}-\frac{7}{8}
\bigg]\cdot
\frac{11}{3}\cdot\frac{2}{11}=\\
&=\bigg[
  \frac{2}{3}+\frac{7}{8}+\frac{1}{3}-\frac{7}{8}
\bigg]\cdot
\frac{2}{3}=\\
&=\bigg[
  \frac{2}{3}+\frac{1}{3}
\bigg]\cdot
\frac{2}{3}=1\cdot\frac{2}{3}=\frac{2}{3}
%\end{soluzione}
\end{flalign*}
\end{esempio}

% \clearpage
Le potenze hanno la precedenza sulle altre operazioni, ma quando ci sono 
anche potenze da calcolare, conviene sempre controllare se è possibile usare 
qualche proprietà. Non solo, ma anche quando non è possibile utilizzare le 
proprietà delle potenze, a volte può essere conveniente non eseguire la 
potenza ma scriverla sotto forma di prodotto:

\begin{esempio}
$\big(\frac{2}{3}\big)^4 \cdot \big(\frac{9}{4}\big)^3  = 
-\frac{2\cdot2\cdot2\cdot2\cdot9\cdot9\cdot9}
      {3\cdot3\cdot3\cdot3\cdot4\cdot4\cdot4} = \frac{9}{4}$
\end{esempio}

In questo caso siamo riusciti a risolvere l'espressione senza eseguire
alcuna moltiplicazione.

\begin{esempio}
Calcola il valore della seguente espressione.
   
\begin{align*}
&\bigg[
  \bigg(\frac{7}{5}-\frac{1}{2}\bigg)^{2}:
  \bigg(\frac{9}{10}\bigg)^{2}-
  \bigg(\frac{2}{15}\bigg)^{4}\cdot
  \bigg(\frac{5}{2}\bigg)^2\cdot
  \bigg(\frac{15}{2}\bigg)^{2}
\bigg]^{2}:
\bigg(\frac{10}{9}\bigg)^{2}-
\bigg(1+\frac{8}{5}+\frac{1}{25}\bigg)=&\\
&=\bigg[
  \bigg(\frac{14 - 5}{10}\bigg)^{2}:
  \bigg(\frac{9}{10}\bigg)^{2}-
  \frac{2\cdot2\cdot2\cdot2}{15\cdot15\cdot15\cdot15}\cdot
  \frac{5\cdot5}{2\cdot2}\cdot
  \frac{15\cdot15}{2\cdot2}
\bigg]^{2}:
\bigg(\frac{10}{9}\bigg)^{2}-
\frac{25 + 40 + 1}{25}=\\
&=\bigg[
  \bigg(\frac{9}{10}\bigg)^{2}:
  \bigg(\frac{9}{10}\bigg)^{2}-
  \frac{1}{9}
\bigg]^{2}\cdot
\bigg(\frac{9}{10}\bigg)^{2}-
\frac{66}{25}=\\
&=\bigg[
  1 -\frac{1}{9}
\bigg]^{2}\cdot
\bigg(\frac{9}{10}\bigg)^{2}-
\frac{66}{25}=\\
&=\bigg[\frac{8}{9}\bigg]^{2}\cdot
\bigg(\frac{9}{10}\bigg)^{2}-
\frac{66}{25}=\\
&=\frac{8\cdot8}{9\cdot9}\cdot
\frac{9\cdot9}{10\cdot10}-
\frac{66}{25}=\frac{16}{25}-\frac{66}{25}=-\frac{50}{25}=-2
\end{align*}

\end{esempio}

Se nell'espressione, oltre alle frazioni, ci sono anche numeri decimali 
limitati o periodici, conviene in un primo passaggio trasformare ogni numero
razionale in frazione e poi eseguire il solito calcolo.

\begin{esempio}
Calcola il valore della seguente espressione.
   
\begin{align*}
&3,5 \cdot 0,4-1,2-0,8\bar{6}\cdot\bigg(
   1,\bar{6}+5,8\bar{3}-5,5-\frac{23}{13}
\bigg)=&\\
&=\frac{35}{10} \cdot \frac{4}{10}-\frac{12}{10}-
 \frac{86-8}{90}\cdot\bigg(
   \frac{16-1}{9}+\frac{583-58}{90}-\frac{55}{10}-\frac{23}{13}
 \bigg)=\\
&=\frac{7}{5}-\frac{6}{5}-
 \frac{78}{90}\cdot\bigg(
   \frac{15}{9}+\frac{525}{90}-\frac{55}{10}-\frac{23}{13}
 \bigg)=\\
&=\frac{1}{5}-
 \frac{13}{15}\cdot\bigg(
   \frac{5}{3}+\frac{35}{6}-\frac{11}{2}-\frac{23}{13}
 \bigg)=\\
&=\frac{1}{5}-
 \frac{13}{15}\cdot\bigg(
   \frac{5}{3}+\frac{35}{6}-\frac{11}{2}-\frac{23}{13}
 \bigg)=\\
&=\frac{1}{5}-
 \frac{13}{15}\cdot
   \frac{130+455-429-138}{2\cdot3\cdot13}=\\
&=\frac{1}{5}-\frac{13}{15}\cdot\frac{18}{2\cdot3\cdot13}=
\frac{1}{5}-\frac{1}{5}=0
\end{align*}

\end{esempio}

% \newpage
\input{\folder reali}
